\section{170~V flyback converter derivation}\label{sec:hv-converter}

The 170~V supply uses an LM51561H flyback controller
(\cref{sec:lm5156h}). This section specifies its circuit requirements and
component values. \Cref{sec:tubes} defines NX1 to NX6.
\Cref{tab:hv-designators} defines the remaining designators.

{\footnotesize
\begin{table}[htbp]
\centering
\caption{170~V flyback converter: component designators}\label{tab:hv-designators}
\begin{tabular}{|l|l|}
\hline
\textbf{Designator} & \textbf{Function} \\
\hline
\des{T}{HV} & Flyback transformer \\

\hline
\des{R}{BALDIG} & Digit-tube anode ballast, one per tube (NX1 to NX4) \\

\hline
\des{R}{BALCOL} & Colon-tube anode ballast, one per tube (NX5 and NX6) \\

\hline
\des{Q}{HVSW} & Primary-side switching FET \\

\hline
\des{R}{HVCS} & Current-sense resistor \\

\hline
\des{R}{HVCSF}, \des{C}{HVCSF} & CS-pin filter \\

\hline
\des{R}{HVD1}, \des{R}{HVD2} & Feedback divider, upper (fixed) \\

\hline
\des{R}{HVD3}, \des{RV}{HVTRIM} & Feedback divider, lower (fixed \& trim) \\

\hline
\des{R}{HVRT} & RT (switching-frequency-set) resistor \\

\hline
\des{R}{HVCOMP}, \des{C}{HVC1}, \des{C}{HVC2} & COMP-pin compensation network \\

\hline
\des{R}{HVGON} & Gate-drive resistor (turn-on) \\

\hline
\des{R}{HVGOFF}, \des{D}{HVGOFF} & Turn-off speedup network (bypasses \des{R}{HVGON} on turn-off only) \\

\hline
\des{D}{HVRECT} & Output rectifier diode \\

\hline
\des{R}{HVCL}, \des{C}{HVCL}, \des{D}{HVCL} & RCD clamp (snubber) network \\

\hline
\des{R}{HVBLEED} & Bleeder resistor \\

\hline
\des{C}{HVSS} & Soft-start capacitor \\

\hline
\des{R}{HVUV1}, \des{R}{HVUV2} & Input UVLO divider \\

\hline
\des{Q}{HVINH}, \des{R}{HVINHPU}, \des{R}{HVINHPD} & Display inhibit shunt and its gate divider \\

\hline
\des{Q}{HVREL}, \des{R}{HVRELPD} & Inhibit release switch and its gate pull-down \\

\hline
\des{R}{HVPG} & PGOOD pull-up resistor \\

\hline
\des{C}{HVVCC} & VCC regulator bypass capacitor \\

\hline
\des{C}{HVPUMP}, \des{R}{HVPUMP} & Bias charge-pump coupling capacitor and series resistor \\

\hline
\des{D}{HVPUMP}, \des{D}{HVRESET} & Charge-pump delivery and reset diodes \\
\hline
\des{D}{HVSTART} & BIAS start-up diode from $V_{\mathrm{BUS}}$ \\
\hline
\des{D}{HVZ} & BIAS clamp zener \\
\hline
\des{C}{HVBIAS} & BIAS reservoir capacitor \\
\hline
\des{C}{HVI1}, \des{C}{HVI2} & Input bulk / high-frequency capacitance \\
\hline
\des{C}{HVO1}, \des{C}{HVO2} & Output capacitance \\

\hline
\end{tabular}
\end{table}
}

\begin{figure}[htbp]
  \centering
  \includegraphics[width=\textwidth]{hv}
  \caption{170~V flyback converter: reference schematic}\label{fig:hv}
\end{figure}

\Cref{fig:hv} shows the connections. It does not show the tubes or ballast
resistors \des{R}{BALDIG} and \des{R}{BALCOL}. \Cref{sec:hv-load} derives these
loads on $V_{\mathrm{HV}}$.

The input $V_{\mathrm{BUS}}$ spans \SIrange{4.75}{12.6}{\volt}
(\cref{sec:pd-classes}). The lower limit is the \SI{5}{\volt} contract at
$-5\,\%$. The upper limit is the \SI{12}{\volt} contract at $+5\,\%$. The
nominal input is \SI{9}{\volt}.

Operation at \SI{4.75}{\volt} adds three requirements. The converter needs the
auxiliary bias supply in \cref{sec:hv-house-pump}. It needs the UVLO threshold
in \cref{sec:hv-uvlo}. It also needs the display inhibit in
\cref{sec:hv-inhibit}. The inhibit uses the negotiated current limit because
the contract voltage does not identify the available current.

\subsection{Preliminaries}\label{sec:hv-prelim}

\subsubsection{Switching frequency}\label{sec:hv-fsw-choice}

In DCM, $L_p f_{\mathrm{sw}} I_{\mathrm{pk}}^2 = 2 P_{\mathrm{in}}$. The
transformer sets $L_p$ (\cref{sec:da2032}). The load sets $P_{\mathrm{in}}$
(\cref{sec:hv-budget}). Thus, the switching frequency sets the peak current.
It also sets the margin below the \SI{3.0}{\ampere} saturation rating. Two
constraints limit the frequency. Both scale with $\sqrt{f_{\mathrm{sw}}}$.

Both $t_{\mathrm{on}}$ and $t_{\mathrm{dem}}$ are proportional to
$f_{\mathrm{sw}}^{-1/2}$. The period is proportional to
$f_{\mathrm{sw}}^{-1}$. Therefore, the conduction fraction increases with
$\sqrt{f_{\mathrm{sw}}}$. At the worst condition in \cref{sec:hv-dcm}, DCM
ends at \SI{185}{\kilo\hertz}. The idle fraction falls to $5\,\%$ at
\SI{168}{\kilo\hertz}.

At low frequency, the permitted current-sense resistance decreases
(\cref{sec:hv-cs-sense}). Core protection requires \des{R}{HVCS} to be at least
\SI{37.8}{\milli\ohm}. The power-delivery limit is
$\SI{93}{\milli\volt}/I_{\mathrm{pk}}^{(L\min)}$. This limit increases with
$\sqrt{f_{\mathrm{sw}}}$. The two limits are equal at
\SI{132}{\kilo\hertz}. The specified \SI{39}{\milli\ohm} resistor meets the
power-delivery limit above \SI{141}{\kilo\hertz}.

The loss model does not add a tighter frequency constraint. The design uses
\SI{150}{\kilo\hertz}. This value is inside both permitted ranges. The implemented frequency
is \SI{149.4}{\kilo\hertz}. At this frequency, the commanded peak current is
\SI{2.18}{\ampere}, or $73\,\%$ of $I_{\mathrm{sat}}$.

\subsubsection{RT resistor}\label{sec:hv-fsw}

For the \SI{150}{\kilo\hertz} target of \cref{sec:hv-fsw-choice}, the RT
equation of \cref{sec:lm5156h} gives



\begin{align*}
  R_T &= \frac{2.21\times10^{10}}{f_{RT}} - 955
       = \frac{2.21\times10^{10}}{1.5\times10^{5}} - 955
       = \SI{146.4}{\kilo\ohm}.
\end{align*}
The nearest standard value is $\des{R}{HVRT} = \SI{147}{\kilo\ohm}$. It gives
\begin{align*}
  f_{RT} &= \frac{2.21\times10^{10}}{147\,000 + 955} = \SI{149.4}{\kilo\hertz}
\end{align*}
All subsequent arithmetic uses \SI{149.4}{\kilo\hertz}.


The DA2032-AL is characterized only at \SI{100}{\kilo\hertz}
(\cref{sec:da2032}). The analysis applies its inductance and saturation data at
\SI{149.4}{\kilo\hertz} (inferred). It does not apply the core-loss data at
this frequency. \Cref{sec:hv-op} includes a core-loss allowance.

\subsubsection{Input UVLO}\label{sec:hv-uvlo}

UVLO must release below the \SI{4.75}{\volt} minimum input. The release voltage
must include the LM51561H threshold tolerance. A \SI{4.75}{\volt} release value
could prevent operation at the specified minimum. The design target is
\SI{4.4}{\volt}.

An external divider sets the UVLO/SYNC thresholds. The typical thresholds are
\SI{1.5}{\volt} rising and \SI{1.45}{\volt} falling
(\cref{sec:lm5156h}). The required divider ratio is:
\begin{align*}
  \frac{\des{R}{HVUV1}+\des{R}{HVUV2}}{\des{R}{HVUV2}}
  &= \frac{4.4}{1.5} = 2.93.
\end{align*}

The divider ratio does not set \des{R}{HVUV2}. A value of tens of kilohms gives
low quiescent current. It also gives useful hysteresis from the
\SI{5}{\micro\ampere} pin current (\cref{sec:lm5156h}). The design uses
\SI{32.4}{\kilo\ohm}:
\begin{align*}
  \des{R}{HVUV1} &= \des{R}{HVUV2}\left(\frac{4.4}{1.5} - 1\right)
  = 32.4 \times 1.933 = \SI{62.6}{\kilo\ohm}
  \;\rightarrow\; \SI{63.4}{\kilo\ohm}.
\end{align*}



The same pin has an EN comparator with an approximately
\SI{0.52}{\volt} rising threshold. The divider refers this threshold to
\SI{1.54}{\volt} at $V_{\mathrm{BUS}}$. Thus, EN becomes true before UVLO
releases during a monotonic input ramp. One divider serves both comparators.

\subsubsection{Display inhibit}\label{sec:hv-inhibit}

UVLO prevents operation before the input rail is present. It does not verify
that the source can supply the tube current. The MCU reads the negotiated
USB-PD contract and controls a display-inhibit signal. N-channel MOSFET
\des{Q}{HVINH} is across \des{R}{HVUV2}. It grounds UVLO/SYNC when it is on. The
MCU releases the inhibit. Two requirements set the gate circuit.

\begin{itemize}[nosep]
  \item \textbf{Assert the inhibit before converter operation.} The MCU and
    converter both use $V_{\mathrm{BUS}}$. However, the MCU 3.3~V rail starts
    after $V_{\mathrm{BUS}}$ appears (\cref{sec:3v3-converter}). A gate bias
    from the 3.3~V rail would depend on two converter start times. A gate bias
    from $V_{\mathrm{BUS}}$ rises with the UVLO divider. Resistor ratios then set
    the sequence.
  \item \textbf{Isolate the MCU control pin from the gate node.} The gate bias
    can reach \SI{11.5}{\volt}. The MCU pin limit is \(V_{DD}+0.3\)~V
    (\cref{sec:esp32}). Thus, the control pin and gate cannot use one node.
\end{itemize}

\des{R}{HVINHPU} and \des{R}{HVINHPD} divide $V_{\mathrm{BUS}}$ onto the gate of
\des{Q}{HVINH}. A second N-channel MOSFET, \des{Q}{HVREL}, pulls this gate to
ground when the MCU enables the display. \des{R}{HVRELPD} holds
\des{Q}{HVREL} off when the MCU pin is high impedance. Thus, a high-impedance
MCU pin keeps the display inhibited. \Cref{sec:esp32-usb} gives the same safe
state to the USB data switch.

\paragraph{Gate divider.} $\des{R}{HVINHPU} = \SI{100}{\kilo\ohm}$ and
$\des{R}{HVINHPD} = \SI{1}{\mega\ohm}$ put $0.909\,V_{\mathrm{BUS}}$ on the
gate nominally. The nominal gate voltage spans \SIrange{4.32}{11.5}{\volt}.
Thus, the required
gate rating is $V_{GS} \geq \pm\SI{12}{\volt}$. The pair draws a maximum of
\SI{11.5}{\micro\ampere}.

\paragraph{Ordering.} The earliest UVLO release includes both divider tolerances
and the minimum rising threshold. It occurs at
$V_{\mathrm{BUS}} = \SI{4.16}{\volt}$. The minimum inhibit-gate divider ratio is
0.900. The gate voltage is then \SI{3.74}{\volt}. Therefore, the shunt requires
$V_{GS(th)} \leq \SI{3.74}{\volt}$. The specified
$V_{GS(th)} \leq \SI{2}{\volt}$ meets this limit. It asserts the inhibit by
$V_{\mathrm{BUS}} = \SI{2.22}{\volt}$. The inhibit therefore becomes active at least
\SI{1.94}{\volt} before UVLO can
release. EN can become true first. Until inhibit assertion, the controller stays
in standby and does not switch (\cref{sec:lm5156h}).

\paragraph{Shunt conduction.} At maximum input, \des{Q}{HVINH} sinks
\SI{199}{\micro\ampere}. It must hold UVLO/SYNC below \SI{1.45}{\volt}. Thus,
its channel resistance must be less than \SI{7.3}{\kilo\ohm}. A part with
single-digit ohm resistance at the \SI{3.74}{\volt} minimum required drive meets this
requirement.

\paragraph{Release switch.} \des{Q}{HVREL} blocks a maximum of
\SI{11.5}{\volt} when off. It sinks a maximum of \SI{126}{\micro\ampere} when
on. The MCU supplies a \SI{3.3}{\volt} gate drive. Therefore,
$V_{GS(th)} \leq \SI{2}{\volt}$. One part can serve both switch positions. A
maximum channel resistance of \SI{4.8}{\kilo\ohm} keeps \des{Q}{HVINH} off.
$\des{R}{HVRELPD} = \SI{100}{\kilo\ohm}$ holds \des{Q}{HVREL} off. The MCU
sinks \SI{33}{\micro\ampere} when it turns \des{Q}{HVREL} on.

\subsubsection{Controller housekeeping (BIAS, VCC, PGOOD)}\label{sec:hv-house}

\paragraph{BIAS.}\label{sec:hv-house-bias}
BIAS supplies the internal regulator for the \SI{6.85}{\volt} VCC gate-driver
rail. The datasheet specifies the VCC output at
$V_{\mathrm{BIAS}} = \SI{8}{\volt}$ (\cref{sec:lm5156h}). A direct connection
to $V_{\mathrm{BUS}}$ cannot maintain this condition below approximately
\SI{8}{\volt}. Therefore, the charge pump in \cref{sec:hv-house-pump} supplies
BIAS. A diode-OR connection from $V_{\mathrm{BUS}}$ supplies start-up current.

\paragraph{VCC regulator capacitor.}\label{sec:hv-house-vcc} The datasheet
requires \SIrange{1}{4.7}{\micro\farad} (\ds{lm5156h.pdf}{14}). The design uses
$\des{C}{HVVCC} = \SI{2.2}{\micro\farad}$. This value allows for the DC-bias
derating of a small ceramic capacitor.

\paragraph{PGOOD pull-up.}\label{sec:hv-house-pgood} PGOOD is an open-drain
output (\cref{sec:lm5156h}). The datasheet recommends a pull-up of at least
\SI{10}{\kilo\ohm}. The pull-up must also charge the pin before the
\SI{25}{\micro\second} deglitch interval ends. The design uses
$\des{R}{HVPG} = \SI{100}{\kilo\ohm}$ to the 3.3~V rail. With
\SI{20}{\pico\farad} of stray capacitance, the time constant is approximately
\SI{2}{\micro\second}. The pull-up voltage meets the PGOOD voltage limits. The
low-state current is \SI{33}{\micro\ampere}, below the
\SI{1}{\milli\ampere} limit.

\paragraph{Spread spectrum (DITHOFF).}\label{sec:hv-house-dithoff}
DITHOFF controls the spread-spectrum function (\ds{lm5156h.pdf}{3}). The
modulation span is $0.156 f_{\mathrm{sw}}$, or approximately
\SI{23}{\kilo\hertz}. It combines a triangular sweep with cycle-to-cycle random
modulation (\ds{lm5156h.pdf}{17}).

The dither decreases the DCM idle time. At the top of the dither range,
$t_{\mathrm{on}}+t_{\mathrm{dem}}$ decreases by $3.7\,\%$. The period decreases
by $7.2\,\%$. Thus, the minimum DCM idle fraction decreases from $10.5\,\%$ to
$7.0\,\%$ (\cref{sec:hv-dcm}). The design does not need spread-spectrum
operation. DITHOFF connects to VCC to disable this function.

\subsubsection{Auxiliary bias charge pump}\label{sec:hv-house-pump}

\paragraph{Requirement.} The gate drive above the Miller plateau turns
\des{Q}{HVSW} off. The plateau and \des{D}{HVGOFF} drop subtract
\SI{3.5}{\volt} from VCC. The remaining voltage drives
$\des{R}{HVGOFF}+R_{\mathrm{int}} = \SI{4.8}{\ohm}$. Thus,
$t_{\mathrm{ovl}} = Q_{GD}/I_{\mathrm{gate}}$. Turn-off overlap is the largest
loss term. The commanded peak current must remain below the
\SI{2.36}{\ampere} minimum current limit (\cref{sec:hv-cs-sense}). Solving the
loss iteration limits overlap to \SI{118}{\nano\second}. This condition
requires at least \SI{4.52}{\volt} of gate drive. A direct BIAS feed from
$V_{\mathrm{BUS}}$ gives only \SI{4.20}{\volt} at minimum input. It cannot meet
the requirement.

Therefore, the switch node supplies BIAS through a charge pump. During the off
time, the node is one tenth of the output voltage above $V_{\mathrm{BUS}}$
(\cref{sec:hv-op-converge}). This level is approximately \SI{24}{\volt} at
minimum input. The pump keeps BIAS above \SI{8}{\volt}. At this voltage, the
datasheet specifies the regulated \SI{6.85}{\volt} VCC output.

\paragraph{Topology.} \des{C}{HVPUMP} and \des{R}{HVPUMP} couple the switch
drain to the pump node. \des{D}{HVPUMP} transfers charge to BIAS.
\des{C}{HVBIAS} stores the charge. \des{D}{HVRESET} returns the pump node to
ground during each on time. \des{D}{HVSTART} supplies BIAS from
$V_{\mathrm{BUS}}$ before switching starts.

During the on time, the drain is near ground. \des{D}{HVRESET} holds the pump
node one diode drop below the drain. This condition charges \des{C}{HVPUMP} to
approximately \SI{0.9}{\volt}. At turn-off, the drain rises to
$V_{\mathrm{BUS}} + V_{\mathrm{refl}}$. The pump node rises with it.
\des{D}{HVPUMP} conducts until the node falls to $V_{\mathrm{BIAS}} + V_F$.

The pump needs no regulation loop, because the controller fixes the charge it
must deliver:
\begin{align*}
  Q_{\mathrm{BIAS}} &= Q_G + \frac{I_Q}{f_{\mathrm{sw}}}
  = 40 + \frac{580\,\si{\micro\ampere}}{149.4\,\si{\kilo\hertz}}
  = \SI{43.9}{\nano\coulomb}
  \quad (I_{\mathrm{BIAS}} = \SI{6.56}{\milli\ampere}),
\end{align*}
so BIAS settles wherever the delivered packet equals that demand:
\begin{align*}
  V_{\mathrm{BIAS}} &= (V_{\mathrm{BUS}} + V_{\mathrm{refl}})
  - \SI{0.9}{\volt} - V_{F} - \frac{Q_{\mathrm{BIAS}}}{k\,\des{C}{HVPUMP}},
  &
  k &= 1 - e^{-t_{\mathrm{dem}}/(\des{R}{HVPUMP}\des{C}{HVPUMP})},
\end{align*}
with $k$ the fraction of the packet transferred inside the demagnetization
window. \des{D}{HVSTART} reverse-biases once the pump lifts BIAS past
$V_{\mathrm{BUS}}$, and must be a Schottky: at \SI{4.75}{\volt} a silicon drop
would leave only \SI{0.55}{\volt} above the controller's \SI{3.5}{\volt} BIAS
minimum during start-up.

\paragraph{Sizing.} \des{R}{HVPUMP}'s time constant must be long against the
\SI{10}{\nano\second} leakage spike (\cref{sec:hv-clamp}) and short against
$t_{\mathrm{dem}}$, so the packet completes inside the shelf. $\des{R}{HVPUMP} =
\SI{82}{\ohm}$ gives $\tau = \SI{385}{\nano\second}$, between that spike and a
$\ge\SI{940}{\nano\second}$ window, and limits the current drawn from the switch
node to \SI{118}{\milli\ampere}.

\des{C}{HVPUMP} is set by the worst delivery corner, the bottom of the input
window at the lowest trim setpoint (\SI{164.8}{\volt},
\cref{tab:hv-ballast-divider}), where the reflected shelf and the load are both
smallest, shortening $t_{\mathrm{dem}}$ to \SI{940}{\nano\second} and $k$ to
\num{0.92}. BIAS must stay at or above \SI{8}{\volt}, the condition at which the
datasheet specifies the \SI{6.85}{\volt} VCC output, and $\des{C}{HVPUMP} =
\SI{4.7}{\nano\farad}$ holds \SI{9.09}{\volt} there even at the $-5\,\%$ limit
of a C0G tolerance. The dielectric is specified because the pump meters charge:
a class-II part's voltage coefficient would change the delivered packet with
operating point.

Across the rest of the window BIAS rises with the input, reaching
\SI{20.4}{\volt} at \SI{12.6}{\volt} in, well inside the pin's \SI{60}{\volt}
recommended range. \des{D}{HVZ}, a \SI{30}{\volt} zener, bounds a fault that
drives the reflected voltage above that range and conducts no current in normal
operation. \des{Q}{HVSW} bounds the fault in turn: the drain reaches its
\SI{100}{\volt} rating at roughly \SI{520}{\volt} out, where the shelf is
\SI{65}{\volt} and the pump meters $k\,\des{C}{HVPUMP}(V_{\mathrm{shelf}} -
\SI{1.6}{\volt} - V_Z) = \SI{149}{\nano\coulomb}$ against the
\SI{43.9}{\nano\coulomb} demand. The zener sinks the \SI{105}{\nano\coulomb}
surplus, $30 \times 105\,\si{\nano\coulomb} \times 149.4\,\si{\kilo\hertz} =
\SI{471}{\milli\watt}$, and a \SI{0.5}{\watt} part covers it as a transient:
charging $C_{\mathrm{out}}$ to \SI{520}{\volt} at the current limit takes about
\SI{140}{\milli\second}.

\des{C}{HVBIAS} carries the controller between packets, at a droop of
$I_{\mathrm{BIAS}}/(\des{C}{HVBIAS} f_{\mathrm{sw}}) = \SI{88}{\milli\volt}$ for
\SI{0.5}{\micro\farad}, negligible against the regulator's \SI{2.2}{\volt} of
headroom at the worst corner, so \SI{0.5}{\micro\farad} \emph{effective} is the
floor. A class-II ceramic at \SI{20.4}{\volt} of bias retains well over half its
value on a \SI{50}{\volt} rating, so \SI{1}{\micro\farad} X7R keeps that floor.
The pump stops during pulse skipping at light load. BIAS then draws current
through \des{D}{HVSTART}. The reduced gate voltage increases turn-off loss, but
the peak current is low.

The pump diodes hold the pump node between $-V_F$ and
$V_{\mathrm{BIAS}}+V_F$. Each diode blocks at most \SI{21.1}{\volt} during
normal operation. \des{D}{HVZ} limits the fault stress to
\SI{30.7}{\volt}. With \SI{20}{\percent} derating, require a
\SI{40}{\volt} rating. Each diode must recover before the next switching edge.

\paragraph{Cost.} Charge enters at the shelf potential and leaves at the
\SI{6.85}{\volt} rail. Thus, the loss is
$Q_{\mathrm{BIAS}}f_{\mathrm{sw}}(V_{\mathrm{shelf}}-\SI{6.85}{\volt})$.
Relative to a direct BIAS feed from $V_{\mathrm{BUS}}$, the increment is:
\begin{align*}
  I_{\mathrm{BIAS}} V_{\mathrm{refl}}
  &= \SI{6.56}{\milli\ampere} \times \SI{19.25}{\volt}
  = \SI{126}{\milli\watt},
\end{align*}
constant across the input window because the shelf tracks $V_{\mathrm{BUS}}$.
\des{C}{HVPUMP} redistributes that loss without changing it: a larger value
settles BIAS higher, moving dissipation out of \des{R}{HVPUMP} and into the
controller's internal regulator. At \SI{12.6}{\volt} in, the split is
\SI{89}{\milli\watt} inside the package and the remainder in \des{R}{HVPUMP}.
Gate charge reduces the overhead because the pump supplies it from the shelf
potential. \Cref{sec:hv-summary} includes this effect.

\subsection{Output voltage}\label{sec:hv-setpoint}

\subsubsection{Output setpoint}\label{sec:hv-ignition}

The IN-12B ignition voltage is at most \SI{170}{\volt}
(\cref{sec:in12b}). This value sets the minimum converter output. The INS-1
limit of \SI{90}{\volt} does not constrain the design. A
\SI{15}{\volt} margin gives an adjustment range of
\SIrange{170}{185}{\volt}. One fixed digit ballast cannot keep the current in
range across this full voltage span (\cref{sec:hv-load-digit}). Therefore,
\cref{sec:hv-divider-split} specifies matched ballast and \des{R}{HVD3} values.
The INS-1 has no rated current range. One ballast value serves all output
settings.

\subsubsection{Interpretation of the datasheet's 200/250~V anode-supply figures}\label{sec:hv-anode-figures}

The IN-12B datasheet specifies a minimum anode supply of \SI{200}{\volt}. It
recommends \SI{250}{\volt} DC. Both values exceed the derived range. These
values can apply to unfiltered half-wave AC (inferred). In that mode, the tube
ignites during each mains cycle and conducts near the voltage peak. The supply
then needs more peak voltage to give sufficient average current.

\subsubsection{Operating points used throughout this section}\label{sec:hv-points}

Two operating points recur, bracketing the trim range:
\begin{description}
  \item[Nominal] $V_{\mathrm{BUS}} = \SI{9}{\volt}$, the intermediate USB-PD
  contract; output at \SI{170}{\volt}, the IN-12B's guaranteed ignition ceiling
  (\cref{sec:hv-ignition}); tube currents at their expected values.
  \item[Design-max envelope] $V_{\mathrm{BUS}} = \SI{4.75}{\volt}$ (bottom of
    the input window, worst case for duty cycle and primary conduction), output
    at \SI{191.5}{\volt}, and every digit tube simultaneously at its
    datasheet-maximum \SI{3.5}{\milli\ampere}, which is what the widest-ballast
    row of \cref{tab:hv-ballast-divider} reaches at exactly that current.
\end{description}

In DCM, the energy packet sets peak current:
$I_{\mathrm{pk}} = \sqrt{2P_{\mathrm{in}}/(L_p f_{\mathrm{sw}})}$. Checks that
depend on input voltage use both limits. The \SI{4.75}{\volt} limit maximizes
duty cycle, conduction loss, and on time. The \SI{12.6}{\volt} limit maximizes
commutation voltage, reverse stress, and sense-current slew rate.

\paragraph{Achieved tuning range.}\label{sec:hv-achieved-range}
Select the digit ballast and feedback resistor together from
\cref{tab:hv-ballast-divider}. Each row maps the rated
\SIrange{2.5}{3.5}{\milli\ampere} current range to a different voltage range.
Thus, $V_{\mathrm{HV}}$ can be set from \SI{164.8}{\volt} to
\SI{191.5}{\volt}. Component stress checks use this full range. Functional
checks use the required \SIrange{170}{185}{\volt} range.

\subsection{Ballast sizing}\label{sec:hv-load}

\subsubsection{Digit tubes (NX1 to NX4, ballast \des{R}{BALDIG})}\label{sec:hv-load-digit}

Each IN-12B anode connects to $V_{\mathrm{HV}}$ through its ballast. The
selected cathode returns through a K155ID1 output. The rated DC cathode current
is $I_K \in \SIrange{2.5}{3.5}{\milli\ampere}$ (\cref{sec:in12b}). The ballast
equation is
\begin{align*}
  R_A &= \frac{V_{\mathrm{rail}} - V_M - V_K}{I_K},
\end{align*}
with maintaining voltage $V_M = \SI{145}{\volt}$ (typical, no tabulated spread;
\cref{sec:in12b}) and $V_K$ the K155ID1's on-state drop at the selected output.
The driver specifies a maximum $V_K$ of \SI{2.5}{\volt} at
\SI{7}{\milli\ampere} (\cref{sec:bcd-driver}). The operating current is less
than half this test current. Therefore, the analysis uses an expected
\SI{1.0}{\volt} drop (inferred). At the smallest ballast, a
\SI{1.5}{\volt} error changes current by approximately
\SI{0.20}{\milli\ampere}. This error is less than the
\SI{0.5}{\milli\ampere} half-width of the rated range.

One ballast value cannot keep the cathode current in its rated range across all
output settings. \Cref{sec:hv-divider-split} therefore gives matched ballast
and \des{R}{HVD3} values. For each row, $V_{\mathrm{min}}$ and
$V_{\mathrm{max}}$ give exactly \SI{2.5}{\milli\ampere} and
\SI{3.5}{\milli\ampere}. The trimmer end stops extend past both targets with
worst-case component tolerances.

The installed \des{R}{BALDIG} value is \SI{7.50}{\kilo\ohm}. The table also
defines permissible alternative values through \SI{13.00}{\kilo\ohm}. Each
alternative requires the matching \des{R}{HVD3}. Calculations that depend on the
design envelope cover the full table. Worst-case dissipation grows with $R_A$ at the shared
\SI{3.5}{\milli\ampere} rated maximum, so the largest entry bounds every row:
\begin{align*}
  I_K^2 R_A = 3.5^2 \times 13.00 \approx \SI{159}{\milli\watt},
\end{align*}
still well inside a DIN0207 axial part's $\approx\SI{0.6}{\watt}$ rating
($\approx$27\,\%).

The IN-12B decimal-point cathode has no driver signal and carries no current.

\subsubsection{Colon tubes (NX5 and NX6, ballast \des{R}{BALCOL})}\label{sec:hv-load-colon}

Each INS-1 anode connects to $V_{\mathrm{HV}}$ through its ballast. Its cathode
connects directly to ground. The colon tubes have no driver. They illuminate
when the high-voltage rail is on. The ballast equation has no driver drop. One
resistor value serves both tubes.

The INS-1 maintaining voltage is at most \SI{55}{\volt}
(\cref{sec:ins1}). Use of this limit gives the minimum calculated current. The
datasheet gives no operating-current range. It gives only a
\SI{0.5}{\milli\ampere} optimum brightness current. The ballast gives this
current at the \SI{177.5}{\volt} center of the output range:
\begin{align*}
  R_A &= \frac{\SI{177.5}{\volt} - \SI{55}{\volt}}{\SI{0.5}{\milli\ampere}}
       = \SI{245}{\kilo\ohm}.
\end{align*}
The nearest standard value is $\des{R}{BALCOL} = \SI{240}{\kilo\ohm}$. Current
spans \SIrange{0.479}{0.542}{\milli\ampere} across the required output range.
It spans \SIrange{0.458}{0.569}{\milli\ampere} across the full adjustment
range. These limits are $8.5\,\%$ below and $14\,\%$ above the optimum.

The datasheet does not give a minimum maintaining voltage. A lower maintaining
voltage increases current through a fixed ballast. For example, a
\SI{45}{\volt} maintaining voltage gives \SI{0.61}{\milli\ampere} at maximum
output. The datasheet gives no maximum current for comparison. Therefore, this
upper current is open.

\paragraph{Ballast tolerance.} \des{R}{BALCOL} has a $5\,\%$ tolerance. This
tolerance gives a current range of \SIrange{0.436}{0.599}{\milli\ampere}. The
datasheet gives no rated range that this variation can exceed. Two resistors at
opposite tolerance limits differ in current by \SI{10.5}{\percent}. A $1\,\%$
pair would differ by $2\,\%$. The datasheet does not relate this current
difference to brightness. A \SI{5}{\volt} difference in maintaining voltage
also causes a $4\,\%$ current difference at midrange. The visible colon
ballasts use carbon-film bodies as a construction preference.

\subsection{Feedback divider design}\label{sec:hv-divider}

\subsubsection{Divider current selection}\label{sec:hv-divider-current}

Two limits set the divider current.

\textbf{Power loss.} The divider loads $V_{\mathrm{HV}}$ whenever the converter
operates. A divider current that is 100 times less than the tube current has a
small power cost.

\textbf{Feedback-node impedance.} A high resistance increases error from FB
bias or leakage current. It also increases coupling from switching noise. The
datasheet gives no FB leakage value. The analysis assumes nanoampere-scale
leakage for this device class (inferred). A \SI{100}{\micro\ampere} divider
current makes this leakage negligible.

At the \SI{170}{\volt} reference point that sets a target total divider
resistance of $(170-1)\,\si{\volt}/\SI{100}{\micro\ampere} \approx
\SI{1.69}{\mega\ohm}$.

\subsubsection{Resistor split and trim range}\label{sec:hv-divider-split}

The divider is $R_{\mathrm{upper}}$ to FB to $R_{\mathrm{lower}}$ to ground.
Its equation is $V_{\mathrm{out}} =
V_{\mathrm{REF}}(1+R_{\mathrm{upper}}/R_{\mathrm{lower}})$. At
\SI{170}{\volt}, the required ratio is 169. The lower resistance is less than
$1\,\%$ of the total. It sets the adjustment range but has little effect on
divider current. Two equal upper resistors share the voltage stress
(\cref{sec:hv-divider-stress}). Therefore, $\des{R}{HVD1} =
\des{R}{HVD2} = \SI{845}{\kilo\ohm}$.

$R_{\mathrm{lower}}$ contains fixed resistor \des{R}{HVD3} and trimmer
\des{RV}{HVTRIM} in series. A \SI{2}{\kilo\ohm} trimmer reaches each row's
$V_{\min}$ and $V_{\max}$ with worst-case divider tolerances. Select
\des{R}{HVD3} so that all three targets remain inside the usable trimmer span.
Adjust the output by measurement.

\begin{table}[htbp]
\centering
\caption{Digit-ballast / \des{R}{HVD3} pairing table (component tolerances: \des{R}{HVD1}/\des{R}{HVD2}
    $\pm1\,\%$ each, \des{R}{HVD3} $\pm1\,\%$, $V_{\mathrm{REF}}$ $\pm1\,\%$,
    \des{RV}{HVTRIM} total resistance assumed $\pm10\,\%$, all stacked
    worst-case additive)}
\label{tab:hv-ballast-divider}
\begin{tabular}{|c|c|c|c|c|c|c|}
\hline
\textbf{$R_A$} & \textbf{$V_{\min}$} & \textbf{$V_{\mathrm{typ}}$} &
\textbf{$V_{\max}$} & \textbf{\des{R}{HVD3}} & \textbf{$V_{\mathrm{floor,wc}}$} &

\textbf{$V_{\mathrm{ceil,wc}}$} \\
(\si{\kilo\ohm}) & (\si{\volt}) & (\si{\volt}) & (\si{\volt}) &
(\si{\kilo\ohm}) & (\si{\volt}) & (\si{\volt}) \\
\hline
7.50  & 164.8 & 168.5 & 172.2 & 9.09 & 146.5 & 192.6 \\
\hline
8.20  & 166.5 & 170.6 & 174.7 & 8.87 & 149.4 & 197.3 \\
\hline
9.10  & 168.8 & 173.3 & 177.8 & 8.87 & 149.4 & 197.3 \\
\hline
10.00 & 171.0 & 176.0 & 181.0 & 8.66 & 152.3 & 202.1 \\
\hline
11.00 & 173.5 & 179.0 & 184.5 & 8.45 & 155.3 & 207.1 \\
\hline
12.00 & 176.0 & 182.0 & 188.0 & 8.25 & 158.3 & 212.1 \\
\hline
13.00 & 178.5 & 185.0 & 191.5 & 8.25 & 158.3 & 212.1 \\
\hline
\end{tabular}
\end{table}

$V_{\min}$, $V_{\mathrm{typ}}$, and $V_{\mathrm{max}}$ are adjustment targets.
They give \SIlist{2.5;3.0;3.5}{\milli\ampere} through the applicable ballast.
Measure the output and adjust it to a target. Divider tolerance does not change
these measured targets. $V_{\mathrm{floor,wc}}$ and
$V_{\mathrm{ceil,wc}}$ are calculated end-stop limits. Start with the
smallest-$R_A$ row. Use a larger ballast only if the installed tubes do not
ignite in that row's range.

\paragraph{User-guide requirement.} The hardware user guide must include
\cref{tab:hv-ballast-divider}. It must identify each row as an acceptable
ballast and \des{R}{HVD3} pair. It must require voltage adjustment with the
tubes removed. The end-stop values are informational. Adjustment to an end stop
is not permitted. The guide must also state that the trimmer range compensates
for component tolerance. The end stops are excluded from operating and stress
analyses.

Divider current at each row's nominal point ranges
\SIrange{99.1}{108.9}{\micro\ampere}, matching the \SI{100}{\micro\ampere}
target of \cref{sec:hv-divider-current} across the whole table.

\subsubsection{Divider element stress}\label{sec:hv-divider-stress}

\des{R}{HVD1} and \des{R}{HVD2} split $(V_{\mathrm{out}}-1)$ evenly:
\SI{84.5}{\volt} each at the \SI{170}{\volt} reference point, \SI{95.25}{\volt}
at the \SI{191.5}{\volt} design-max envelope (\cref{sec:hv-achieved-range}),
against the \SI{250}{\volt} limiting element voltage of the specified DIN0207
axial metal film parts. Dissipation is $95.25^2/\SI{845}{\kilo\ohm} =
\SI{10.7}{\milli\watt}$. \des{R}{HVD3} and \des{RV}{HVTRIM} see about
\SI{1}{\volt}.

\subsection{Aggregate current and power budget}\label{sec:hv-budget}

The bleeder \des{R}{HVBLEED} completes the load list at \SI{1.5}{\mega\ohm},
sized in \cref{sec:hv-bleeder} against the output-discharge and minimum-load
requirements.

{\footnotesize
\begin{table}[htbp]
\centering
\caption{$V_{\mathrm{HV}}$: aggregate anode current budget}\label{tab:hv-current-budget}
\begin{tabular}{|l|r|r|}
\hline
\textbf{Load element} & \textbf{Nominal (\SI{170}{\volt})} & \textbf{Design-max (\SI{191.5}{\volt})} \\
\hline
NX1 to NX4 digit cathodes & $4 \times 3.20 = \SI{12.80}{\milli\ampere}$
  & $4 \times 3.50 = \SI{14.00}{\milli\ampere}$ \\
\hline
NX5 and NX6 colon anodes & $2 \times 0.479 = \SI{0.96}{\milli\ampere}$
  & $2 \times 0.569 = \SI{1.14}{\milli\ampere}$ \\
\hline
Bleeder \des{R}{HVBLEED} (\SI{1.5}{\mega\ohm}) & \SI{0.11}{\milli\ampere} & \SI{0.13}{\milli\ampere} \\

\hline
Feedback divider ($\approx\SI{1.70}{\mega\ohm}$) & \SI{0.10}{\milli\ampere} & \SI{0.11}{\milli\ampere} \\
\hline
Total $I_{\mathrm{out}}$ & \SI{13.97}{\milli\ampere} & \SI{15.38}{\milli\ampere} \\
\hline
$P_{\mathrm{out}} = V_{\mathrm{out}} I_{\mathrm{out}}$ & \SI{2.38}{\watt} & \SI{2.94}{\watt} \\
\hline
\end{tabular}
\end{table}
}

The nominal-point digit and colon currents follow from the specified ballast
values (\cref{sec:hv-load-digit,sec:hv-load-colon}) at \SI{170}{\volt}; the
divider's \SI{0.10}{\milli\ampere} is under 1\,\% of total output current at
either operating point, as \cref{sec:hv-divider-current} assumed.

\textbf{The maximum-design point can occur.} With a
\SI{13.00}{\kilo\ohm} ballast and a \SI{191.5}{\volt} output, each digit tube
draws \SI{3.5}{\milli\ampere}. Thus, all four tubes can operate at their rated
maximum current at the same time. The INS-1 has no rated maximum current
(\cref{sec:ins1}).

\subsection{Operating point and efficiency}\label{sec:hv-op}

\subsubsection{Component values for the loss model}\label{sec:hv-op-inputs}

The loss model uses the following values. Their cross-referenced sections
derive them from the converged peak current.
\begin{itemize}[nosep]
  \item \des{Q}{HVSW} is an IRL540N-type commodity through-hole logic-level power MOSFET;
  \item $\des{R}{HVGON} = \SI{10}{\ohm}$, governing turn-on;
  \item \des{D}{HVGOFF} (SB240S-E3/54) and $\des{R}{HVGOFF} = \SI{3.3}{\ohm}$,
    the turn-off speedup network, sized in
    \cref{sec:hv-op-losses,sec:hv-switch};
  \item $\des{R}{HVCS} = \SI{39}{\milli\ohm}$, the current-sense resistor,
    derived in \cref{sec:hv-cs};
  \item $\des{R}{HVCL} = \SI{30}{\kilo\ohm}$, $\des{C}{HVCL} =
    \SI{12}{\nano\farad}$, the RCD clamp network, sized in
    \cref{sec:hv-clamp}.
\end{itemize}

\subsubsection{Loss mechanisms}\label{sec:hv-op-losses}

In DCM, the converter transfers a fixed magnetizing-energy packet in each
cycle:
\begin{align*}
  P_{\mathrm{in}} = \tfrac{1}{2} L_p I_{\mathrm{pk}}^2 f_{\mathrm{sw}}
  \quad\Longrightarrow\quad
  I_{\mathrm{pk}} = \sqrt{\frac{2 P_{\mathrm{in}}}{L_p f_{\mathrm{sw}}}},
  \qquad
  P_{\mathrm{in}} = \frac{P_{\mathrm{out}}}{\eta}.
\end{align*}
Calculate efficiency by iteration. Start with an efficiency value. Calculate
$P_{\mathrm{in}}$ and $I_{\mathrm{pk}}$. Then calculate each loss and update
the efficiency. Repeat these calculations until they converge. The loss models
are as follows.
\begin{description}[style=nextline]
  \item[Primary conduction] Triangular current over $t_{\mathrm{on}}$:
    $I_{\mathrm{rms}} = I_{\mathrm{pk}}\sqrt{D/3}$ through
    $R_{DS(\mathrm{on})} + \des{R}{HVCS} + R_{\mathrm{DCR,pri}}
     = 55 + 39 + 13 = \SI{107}{\milli\ohm}$. The datasheet tabulates
    $R_{DS(\mathrm{on})}$ at only two $V_{GS}$ points, both at
    $T_J=\SI{25}{\celsius}$: \SI{53}{\milli\ohm} at $V_{GS}=\SI{5.0}{\volt}$ and
    \SI{44}{\milli\ohm} at $V_{GS}=\SI{10}{\volt}$
    (\ds{infineon-irl540n-datasheet-en.pdf}{2}), so interpolating to the
    \SI{6.85}{\volt} gate rail gives a cold value of
    $\approx\SI{49.7}{\milli\ohm}$. It then rises with junction temperature
    (\ds{infineon-irl540n-datasheet-en.pdf}{3}, Fig.\,4, read as a linear
    approximation between its $1.0\times$-at-\SI{25}{\celsius} and
    $2.4\times$-at-\SI{175}{\celsius} anchors), which \des{Q}{HVSW}'s own
    dissipation sets. Solved self-consistently against the design-max loss
    budget below (\SI{25}{\celsius} ambient assumed), the converged point is
    $T_J \approx \SI{37}{\celsius}$, $R_{DS(\mathrm{on})} \approx
    \SI{55}{\milli\ohm}$.

  \item[Turn-off overlap]
    The IRL540N datasheet (\ds{infineon-irl540n-datasheet-en.pdf}{2}) specifies
    the gate-drain (Miller) charge only as a guaranteed maximum, $Q_{GD} \le
    \SI{38}{\nano\coulomb}$, with no typical value tabulated; its gate-charge
    curve (\ds{infineon-irl540n-datasheet-en.pdf}{4}, Fig.\,6) shows the Miller
    plateau sitting at $V_{\mathrm{plateau}} \approx \SI{3}{\volt}$. The same
    datasheet's resistive-load switching-time table gives $t_f =
    \SI{62}{\nano\second}$ typical at its own test condition
    ($R_G=\SI{5.0}{\ohm}$, $V_{GS}$ pulse $=\SI{5.0}{\volt}$), which implies a
    typical $Q_{GD} \approx t_f \times (5.0-3.0)/5.0\,\si{\ohm} \approx
    \SI{25}{\nano\coulomb}$ at that test point. The LM51561H's output stage adds
    an internal pull-down resistance in series with $\des{R}{HVGON}$. Its datasheet
    (\ds{lm5156h.pdf}{6}, MOSFET DRIVER) gives the low-state output's voltage
    drop under an externally forced \SI{100}{\milli\ampere} as \SI{0.15}{\volt},
    implying $R_{\mathrm{int}} \approx \SI{1.5}{\ohm}$ for this pull-down path
    (the table's two rows are ambiguous about which pin state pairs with
    which current direction, so this is the best available reading).
    Turn-off overlap dominates the loss budget, so that edge is given a faster
    discharge path of its own. \des{D}{HVGOFF} in series with
    \des{R}{HVGOFF} bypasses $\des{R}{HVGON}$ for the discharge direction only:
    its anode sits at the gate, its cathode at the
    $\des{R}{HVGON}$/driver-output node, so it blocks during turn-on, leaving
    that edge unchanged, and conducts during turn-off. \des{R}{HVGOFF} is sized
    against the driver's \SI{1.5}{\ampere} peak source/sink capability
    (\ds{lm5156h.pdf}{29}), the binding current being the brief pre-plateau
    discharge from the fully-enhanced gate, which exceeds the plateau current:
    \begin{align*}
      I_{\mathrm{gate}}^{\mathrm{peak}} &= \frac{V_{\mathrm{drive}} -
        V_{F(\des{D}{HVGOFF})}}{\des{R}{HVGOFF} + R_{\mathrm{int}}}
        = \frac{6.85 - 0.5}{3.3+1.5} \approx \SI{1.32}{\ampere}
        \quad (88\,\% \text{ of the } \SI{1.5}{\ampere} \text{ rating}),
    \end{align*}
    using $V_{F(\des{D}{HVGOFF})} \approx \SI{0.5}{\volt}$, interpolated below
    the sole tabulated point ($\SI{0.55}{\volt}$ at \SI{2}{\ampere}). The
    \SI{12}{\percent} headroom rests on that interpolation and on the
    $R_{\mathrm{int}}$ estimate; both are derived, neither measured.
    At the lower Miller-plateau voltage, the turn-off gate current is
    $I_{\mathrm{gate}} = (6.85-3.0-0.5)/(3.3+1.5) \approx \SI{698}{\milli\ampere}$,
    giving
    \begin{align*}
      t_{\mathrm{ovl}}^{\mathrm{typ}} &= \frac{Q_{GD}^{\mathrm{typ}}}{I_{\mathrm{gate}}}
        \approx \SI{36}{\nano\second}, &
      t_{\mathrm{ovl}}^{\max} &= \frac{Q_{GD}^{\max}}{I_{\mathrm{gate}}}
        \approx \SI{54}{\nano\second}.
    \end{align*}
    With the commutation voltage $V_{\mathrm{com}} = V_{\mathrm{BUS}} +
    V_{\mathrm{refl}}$, $E_{\mathrm{off}} \approx \tfrac{1}{2} I_{\mathrm{pk}}
    V_{\mathrm{com}} t_{\mathrm{ovl}}$. The convergence below uses
    $t_{\mathrm{ovl}}^{\mathrm{typ}}$; \cref{sec:hv-switch} checks
    \des{Q}{HVSW}'s dissipation at $t_{\mathrm{ovl}}^{\max}$. Turn-on is
    near-lossless in DCM (zero current; $C_{oss}$ discharge contributes a few
    milliwatts) and is governed by $\des{R}{HVGON}$ alone.
    \item[Gate drive and controller]
    $P = (Q_G f_{\mathrm{sw}} + I_Q)(V_{\mathrm{BUS}} + V_{\mathrm{refl}})$
    with $Q_G \approx \SI{40}{\nano\coulomb}$ at the \SI{6.85}{\volt} VCC rail
    (scaled from the \SI{10}{\volt} datasheet point, inferred) and controller
    operating current $I_Q \le \SI{580}{\micro\ampere}$ (\ds{lm5156h.pdf}{5}). Gate
    charge is drawn through the internal VCC regulator from BIAS, which the
    charge pump of \cref{sec:hv-house-pump} lifts from the switch node, so the
    charge is drawn at the reflected shelf, well above $V_{\mathrm{BUS}}$. It is
    the largest single term in the budget below.

  \item[RCD clamp] Leakage energy $\tfrac{1}{2} L_{\mathrm{lk}} I_{\mathrm{pk}}^2$
    per cycle, dissipated in \des{R}{HVCL} with the clamp-voltage multiplier
    derived in \cref{sec:hv-clamp}.
  \item[Rectifier] $V_F \approx \SI{1.0}{\volt}$ (ultrafast-diode class figure at
    $\sim$\SI{0.2}{\ampere}) times $I_{\mathrm{out}}$; recovery loss is negligible
    because in DCM the secondary current reaches zero before the switch turns on.
  \item[Secondary conduction] $I_{s,\mathrm{rms}}^2 \times \SI{1.60}{\ohm}$ DCR.
  \item[Capacitor ESR] Input $\approx\SI{40}{\milli\ohm}$ and output
  $\approx\SI{0.5}{\ohm}$, both the polymer parts' own; \des{C}{HVI2}'s
  impedance stays above the input figure below \SI{850}{\kilo\hertz}, past the
  ripple's dominant harmonics.
  \item[Core] Uncharacterized at this frequency; a \SI{50}{\milli\watt} allowance
    is carried.
\end{description}

\subsubsection{Converged operating point (design-max envelope)}\label{sec:hv-op-converge}

Iterating at the design-max envelope ($P_{\mathrm{out}} = \SI{2.94}{\watt}$,
$V_{\mathrm{BUS}} = \SI{4.75}{\volt}$, $V_{\mathrm{out}} = \SI{191.5}{\volt}$,
reflected voltage $V_{\mathrm{refl}} = (191.5+1)/10 = \SI{19.25}{\volt}$,
$t_{\mathrm{ovl}} = t_{\mathrm{ovl}}^{\mathrm{typ}} = \SI{36}{\nano\second}$),
the loop converges to $\eta = \SI{83.3}{\percent}$:
\begin{align*}
  P_{\mathrm{in}} &= \frac{2.945}{0.833} = \SI{3.54}{\watt}, &
  E_{\mathrm{cyc}} &= \frac{3.54}{149.4\times10^{3}} = \SI{23.7}{\micro\joule},\\
  I_{\mathrm{pk}} &= \sqrt{\frac{2 \times 23.7\,\si{\micro\joule}}
  {\SI{10}{\micro\henry}}} = \SI{2.18}{\ampere}, &
  t_{\mathrm{on}} &= \frac{L_p I_{\mathrm{pk}}}{V_{\mathrm{BUS}}}
                   = \frac{10\,\si{\micro\henry} \times 2.18}{4.75}
                   = \SI{4.58}{\micro\second},\\
                   D &= t_{\mathrm{on}} f_{\mathrm{sw}} = 0.68, &
                   I_{\mathrm{rms,pri}} &= I_{\mathrm{pk}}\sqrt{D/3} = \SI{1.04}{\ampere},\\
                   I_{\mathrm{in}} &= \frac{P_{\mathrm{in}}}{V_{\mathrm{BUS}}}
                   = \SI{744}{\milli\ampere}.
\end{align*}
Converged loss budget at this point:

{\footnotesize
\begin{table}[htbp]
\centering
\caption{170~V flyback converter: loss budget at the design-max envelope}\label{tab:hv-loss-budget}
\begin{tabular}{|l|r|}
\hline
\textbf{Mechanism} & \textbf{Loss} \\
\hline
Bias draw through the charge pump (\cref{sec:hv-house-pump}) & \SI{157}{\milli\watt} \\

\hline
MOSFET turn-off overlap ($t_{\mathrm{ovl}}=\SI{36}{\nano\second}$) & \SI{140}{\milli\watt} \\

\hline
Primary conduction (FET 60, \des{R}{HVCS} 42, DCR 14) & \SI{116}{\milli\watt} \\

\hline
RCD clamp (\des{R}{HVCL}) & \SI{86}{\milli\watt} \\

\hline
Core allowance (uncharacterized) & \SI{50}{\milli\watt} \\
\hline
Capacitor ESR (input + output) & \SI{22}{\milli\watt} \\

\hline
Rectifier \des{D}{HVRECT} conduction & \SI{15}{\milli\watt} \\

\hline
Secondary DCR & \SI{4}{\milli\watt} \\
\hline
Total & \SI{590}{\milli\watt} \\

\hline
$\eta = 2.945/(2.945+0.590)$ & \SI{83.3}{\percent} \\

\hline
\end{tabular}
\end{table}
}

\paragraph{Efficiency curve.}\label{sec:hv-op-curve}
The charge pump holds the gate rail constant (\cref{sec:hv-house-pump}).
Turn-off overlap and bias loss increase with
$V_{\mathrm{BUS}} + V_{\mathrm{refl}}$. Together, they increase from
\SI{297}{\milli\watt} to \SI{400}{\milli\watt}. Primary conduction and
capacitor ESR decrease with duty cycle. Together, they decrease from
\SI{138}{\milli\watt} to \SI{57}{\milli\watt}. The other modeled losses total
\SI{155}{\milli\watt}. The total loss is
$A(V_{\mathrm{BUS}}+V_{\mathrm{refl}})+B/V_{\mathrm{BUS}}+C$. Here,
$A=\SI{12.4}{\milli\watt\per\volt}$ and
$B=\SI{683}{\milli\watt\volt}$. The minimum occurs at
$\sqrt{B/A}=\SI{7.4}{\volt}$:

{\footnotesize
  \begin{table}[htbp]
    \centering
    \caption{Design-max efficiency across the input window}\label{tab:hv-eta-line}
    \begin{tabular}{|l|c|c|c|c|c|}
      \hline
      $V_{\mathrm{BUS}}$ & \SI{4.75}{\volt} & \SI{7}{\volt} & \SI{9}{\volt} &
      \SI{11}{\volt} & \SI{12.6}{\volt} \\
      \hline
      Duty & 0.68 & 0.46 & 0.36 & 0.30 & 0.26 \\
      \hline
      Total loss & \SI{590}{\milli\watt} & \SI{578}{\milli\watt} &
      \SI{583}{\milli\watt} & \SI{596}{\milli\watt} & \SI{610}{\milli\watt} \\
      \hline
      $\eta$ & \SI{83.3}{\percent} & \SI{83.6}{\percent} & \SI{83.5}{\percent} &
      \SI{83.2}{\percent} & \SI{82.8}{\percent} \\
      \hline
    \end{tabular}
  \end{table}
}

The convergence above uses $t_{\mathrm{ovl}}^{\mathrm{typ}}$; across the
\SIrange{36}{54}{\nano\second} spread the design-max efficiency band runs
\SIrange{81.4}{83.3}{\percent}. \Cref{sec:hv-switch} checks \des{Q}{HVSW}'s
dissipation at $t_{\mathrm{ovl}}^{\max}$, and \cref{sec:hv-cs-sense} folds the
same spread into its current-limit margin as a tolerance term.

The two largest losses depend on the \des{Q}{HVSW} gate charge. Bias loss
depends on $Q_G$. Turn-off overlap depends on $Q_{GD}$. Together, these losses
are \SI{297}{\milli\watt} of the \SI{590}{\milli\watt} total. Both values are
derived. Measure converter efficiency to verify this loss model.

\subsubsection{Nominal-point efficiency}\label{sec:hv-op-nominal}

The same procedure at the nominal point ($\SI{2.38}{\watt}$, \SI{170}{\volt})
gives, at each contract:

{\footnotesize
  \begin{table}[htbp]
    \centering
    \caption{Nominal-point ($\SI{2.38}{\watt}$, \SI{170}{\volt}) operating figures}\label{tab:hv-nominal-line}
    \begin{tabular}{|l|c|c|c|}
      \hline
      $V_{\mathrm{BUS}}$ & \SI{4.75}{\volt} & \SI{9}{\volt} & \SI{12.6}{\volt} \\
      \hline
      $\eta$ & \SI{82.7}{\percent} & \SI{82.5}{\percent} & \SI{81.7}{\percent} \\
      \hline
      $P_{\mathrm{in}}$ & \SI{2.87}{\watt} & \SI{2.88}{\watt} & \SI{2.91}{\watt} \\
      \hline
      $I_{\mathrm{pk}}$ & \SI{1.96}{\ampere} & \SI{1.96}{\ampere} & \SI{1.97}{\ampere} \\
      \hline
      Duty & 0.62 & 0.33 & 0.23 \\
      \hline
      $I_{\mathrm{rms,pri}}$ & \SI{0.89}{\ampere} & \SI{0.65}{\ampere} & \SI{0.55}{\ampere} \\
      \hline
      $I_{\mathrm{in}}$ average & \SI{605}{\milli\ampere} & \SI{320}{\milli\ampere} &
      \SI{231}{\milli\ampere} \\
      \hline
    \end{tabular}
  \end{table}
}

Nominal efficiency is lower than maximum-load efficiency. The main losses are
proportional to $\sqrt{P}$. Bias current does not decrease with load. Therefore,
the \SIrange{143}{195}{\milli\watt} bias loss is a larger fraction of the
\SI{2.38}{\watt} nominal output. The nominal efficiency is
0.6 to 1.1 percentage points lower than the maximum-load efficiency.

The $I_{\mathrm{in}}$ row is what \cref{sec:hv-inhibit}'s current test is against.
A \SI{5}{\volt} contract must carry \SI{605}{\milli\ampere} at the nominal point
and \SI{744}{\milli\ampere} at the design-max envelope, so a
\SI{3}{\ampere} \SI{5}{\volt} contract qualifies and a default
\SI{500}{\milli\ampere} port does not.

\subsubsection{Peak current at minimum inductance}\label{sec:hv-op-lmin}

With the minimum-tolerance inductance ($L_p = \SI{9}{\micro\henry}$ at rated peak
current, \cref{sec:da2032}), the same energy packet requires
\begin{align*}
  I_{\mathrm{pk}}^{(L\min)} &= \sqrt{\frac{2 \times 23.7\,\si{\micro\joule}}
                                          {\SI{9}{\micro\henry}}}
                             = \SI{2.29}{\ampere},
\end{align*}
the worst-case commanded peak used in \cref{sec:hv-cs}.

\subsection{DCM verification}\label{sec:hv-dcm}

DCM requires $t_{\mathrm{on}} + t_{\mathrm{dem}} < T$. The remaining time must
include the drain ring-down. For the 1:10 turns ratio, the demagnetization time
is
\begin{align*}
  t_{\mathrm{dem}} &= \frac{n L_p I_{\mathrm{pk}}}{V_{\mathrm{out}} + V_F}.
\end{align*}
The binding corner maximizes $t_{\mathrm{on}} + t_{\mathrm{dem}}$: maximum
inductance ($\SI{11}{\micro\henry}$, $+10\,\%$), minimum input
(\SI{4.75}{\volt}), design-max load:
\begin{align*}
  I_{\mathrm{pk}} &= \sqrt{\frac{2\times23.7\,\si{\micro\joule}}{\SI{11}{\micro\henry}}}
  = \SI{2.07}{\ampere},\\
  t_{\mathrm{on}} &= \frac{11\,\si{\micro\henry}\times 2.07}{4.75}
  = \SI{4.80}{\micro\second}, &
  t_{\mathrm{dem}} &= \frac{10 \times 11\,\si{\micro\henry} \times 2.07}{192.5}
  = \SI{1.19}{\micro\second},\\
  t_{\mathrm{on}} + t_{\mathrm{dem}} &= \SI{5.99}{\micro\second}
  = 0.89\,T, &
  t_{\mathrm{idle}} &= \SI{0.71}{\micro\second}.
\end{align*}

\textbf{DCM has a minimum idle fraction of $10.5\,\%$.} The secondary circuit
sets demagnetization time, which does not depend on input voltage. On time is
inversely proportional to $V_{\mathrm{BUS}}$. Therefore, minimum input gives
the smallest idle fraction. At \SI{12.6}{\volt}, duty falls to 0.26. A lighter
load also increases the idle fraction.

CCM operation invalidates the energy-packet loss model and the single-pole
plant in \cref{sec:hv-comp}. It also invalidates the rectifier-recovery and
slope-compensation conclusions. The \SI{0.71}{\micro\second} idle time can end
during drain ring-down. This timing adds only a few milliwatts because DCM
turn-on current is zero.

The 0.72 duty cycle does not cause subharmonic oscillation in DCM. The
$D>0.5$ threshold applies to CCM. In DCM, the current returns to zero during
each cycle. Therefore, the design does not require $R_{SL}$
(\cref{sec:hv-cs}).

Duty at the worst corner remains under the controller's $\approx$0.91 maximum
near this frequency, and $t_{\mathrm{on}} \ge \SI{4.8}{\micro\second}$ is well
above the $\approx\SI{165}{\nano\second}$ minimum on-time (\cref{sec:lm5156h}).

\subsection{Current sensing and current limit}\label{sec:hv-cs}

\subsubsection{CS filter}\label{sec:hv-cs-filter}

With the \SI{100}{\milli\volt} (93 to 107) current-limit threshold of
\cref{sec:lm5156h} and \des{R}{HVCS} $= \SI{39}{\milli\ohm}$ (declared in
\cref{sec:hv-op-inputs}):
\begin{align*}
  I_{\mathrm{lim}}^{\mathrm{typ}} &= \frac{\SI{100}{\milli\volt}}{\SI{39}{\milli\ohm}}
    = \SI{2.56}{\ampere},\\
  I_{\mathrm{lim}}^{\min} &= \frac{\SI{93}{\milli\volt}}{39\,\si{\milli\ohm}\times1.01}
    = \SI{2.36}{\ampere}, &
  I_{\mathrm{lim}}^{\max} &= \frac{\SI{107}{\milli\volt}}{39\,\si{\milli\ohm}\times0.99}
    = \SI{2.77}{\ampere}.
\end{align*}
The CS pin connects to the sense node through \des{R}{HVCSF}.
\des{C}{HVCSF} connects the pin to ground. The circuit does not use $R_{SL}$.
Thus, the current-limit threshold does not depend on duty cycle. The minimum
on-time and off-time criteria have large margins. Two other requirements set
$\tau = \des{R}{HVCSF}\des{C}{HVCSF}$.

\textbf{Upper limit: current-limit error.} The filter delays the sampled current
by $\Delta I_{\mathrm{delay}} = (V_{\mathrm{BUS}}/L_p)\,\tau$. Current continues
to increase during this delay. Limit the increase to half of the current-limit
tolerance range, or \SI{0.205}{\ampere}:
\begin{align*}
  \tau < \frac{\SI{0.205}{\ampere}}{(V_{\mathrm{BUS}}/L_p)_{\max}}
       = \frac{\SI{0.205}{\ampere}}{\SI{1.40}{\ampere\per\micro\second}}
       = \SI{146}{\nano\second},
\end{align*}
evaluated at the fastest observed slope (\SI{12.6}{\volt}, \SI{9}{\micro\henry}).

\textbf{Lower limit: spike attenuation.} Gate current and $C_{oss}$ discharge
cause a turn-on current spike. The spike duration is tens of nanoseconds
(inferred). The filter time constant must be comparable to this duration.

$\des{R}{HVCSF} = \SI{100}{\ohm}$ and
$\des{C}{HVCSF} = \SI{470}{\pico\farad}$. Both values are in the TI
recommended ranges. Their time constant is
\begin{align*}
  \tau = \des{R}{HVCSF}\des{C}{HVCSF} = 100 \times 470\,\si{\pico\farad} = \SI{47}{\nano\second},
\end{align*}
This value is comparable to the inferred spike duration. It is less than the
\SI{146}{\nano\second} maximum. The resulting current error is approximately
\SI{66}{\milli\ampere}. A \SI{50}{\nano\second} comparator and driver delay
(inferred) adds \SI{70}{\milli\ampere}. The total overshoot is approximately
\SI{0.14}{\ampere}.
\Cref{sec:hv-cs-sense} uses this value.

\subsubsection{Sense resistor and limit thresholds}\label{sec:hv-cs-sense}

\des{R}{HVCS} carries the switch current to ground; with the current-limit
thresholds of \cref{sec:hv-cs-filter} and the filter's overshoot $\Delta I_{\mathrm{ovs}}
\approx \SI{0.14}{\ampere}$, the sense resistor must satisfy two opposing
constraints:
\begin{align*}
  \text{power delivery:}\quad
    R_{S} &< \frac{\SI{93}{\milli\volt}}{1.01 I_{\mathrm{pk}}^{(L\min)}}
          = \frac{93\,\si{\milli\volt}}{1.01\times2.29\,\si{\ampere}} = \SI{40.2}{\milli\ohm},\\
  \text{core protection:}\quad
    R_{S} &> \frac{\SI{107}{\milli\volt}}{0.99(I_{\mathrm{sat}} - \Delta I_{\mathrm{ovs}})}
          = \frac{107\,\si{\milli\volt}}{0.99(3.0 - 0.14)\,\si{\ampere}} = \SI{37.8}{\milli\ohm},
\end{align*}
a \SI{2.4}{\milli\ohm} nominal window; neither bound depends on
$V_{\mathrm{out}}$. \SI{39}{\milli\ohm} (declared in \cref{sec:hv-op-inputs})
sits inside it, \SI{1.2}{\milli\ohm} under the power-delivery ceiling and
\SI{1.2}{\milli\ohm} above the core-protection floor. Both bounds are
worst-case-additive: minimum datasheet threshold, \des{R}{HVCS} at its own
tolerance extreme, and, on the power-delivery side, $L_p$ at its $-10\,\%$
extreme, all stacked simultaneously in the same direction.

\SI{39}{\milli\ohm} is the only standard value in this range.

\textbf{Tolerance combination.} An insufficient margin causes the current
limit to end a switching cycle early. The result is reduced output energy and
display brightness. It does not cause component damage. Therefore, this check
also uses a root-sum-square (RSS) tolerance combination. Turn-off overlap is an
independent part-to-part variation. Include it in the RSS calculation.

The nominal margin is $M_{\mathrm{nom}}=\SI{0.389}{\ampere}$. It uses
$\SI{100}{\milli\volt}/\SI{39}{\milli\ohm}$ and nominal $L_p$. It also uses
the typical turn-off overlap. The individual changes are:
\begin{align*}
  \Delta_{\mathrm{threshold}} &= \SI{-0.180}{\ampere}, &
  \Delta_{R_S} &= \SI{-0.025}{\ampere}, &
  \Delta_{L_p} &= \SI{-0.118}{\ampere}, &
  \Delta_{t_{\mathrm{ovl}}} &= \SI{-0.026}{\ampere}.
\end{align*}
Worst-case-additive sums these directly: $M_{\mathrm{wc}} = 0.389 - 0.180 - 0.025
- 0.118 - 0.026 = \SI{+0.040}{\ampere}$ ($+1.7\,\%$ of $I_{\mathrm{lim,min}}$).
RSS combines the same four in quadrature:
\begin{align*}
  M_{\mathrm{RSS}} &= M_{\mathrm{nom}} - \sqrt{\Delta_{\mathrm{threshold}}^2 +
    \Delta_{R_S}^2 + \Delta_{L_p}^2 + \Delta_{t_{\mathrm{ovl}}}^2}
    = 0.389 - 0.218 = \SI{+0.171}{\ampere}
    \quad (+7.2\,\% \text{ of } I_{\mathrm{lim,min}}),
\end{align*}
The threshold specification dominates the combination. \textbf{The RSS treatment
is scoped to this one check}; every voltage- and thermal-rating margin elsewhere
remains worst-case-additive, because those checks guard against destructive
outcomes.

Both tolerance methods give a positive margin. The additive margin is
$1.7\,\%$. The RSS margin is $7.2\,\%$. The bias current decreases this margin
because it increases input power and commanded peak current.

\paragraph{What hiccup changes about that.}\label{sec:hv-cs-hiccup}
On the LM51561H, sixty-four \emph{consecutive} limited cycles trip the fault
timer (\cref{sec:hv-ss}), after which the converter stops for
\SI{219}{\milli\second} and retries. A board that exhausts this margin at the
top of its trim range on a \SI{5}{\volt} contract blinks at roughly
\SI{4.6}{\hertz}.

The fault requires 64 consecutive limited cycles. Eight unlimited cycles reset
the counter. Thus, intermittent limiting does not cause hiccup. Continuous
limiting requires the maximum load and output setting. A lower output setting
or a larger ballast decreases peak current.

\textbf{The hardware user guide must include this condition.} If the display
blinks at maximum adjustment on a \SI{5}{\volt} supply, decrease the output.
Alternatively, install a larger ballast. Measure current limit with the
selected transformer to verify the calculated margin.

On the other side, the maximum clamped peak plus overshoot reaches
$2.77 + 0.14 = \SI{2.91}{\ampere}$ against the \SI{3.0}{\ampere} transformer
rating: a 3\,\% margin, acceptable because the \SI{3.0}{\ampere} figure is a
soft-saturation point, with $L$ already specified at \SI{9}{\micro\henry}
there.

\des{R}{HVCS} dissipation is $I_{\mathrm{rms}}^2 R_S = 1.04^2 \times 0.039 =
\SI{42}{\milli\watt}$ at the design-max envelope, trivial in a 2512.

\subsection{RCD clamp}\label{sec:hv-clamp}

At turn-off, the leakage inductance carries $I_{\mathrm{pk}}$. This inductance
is not coupled to the secondary. \des{D}{HVCL} directs its current into
\des{C}{HVCL}. \des{R}{HVCL} discharges the capacitor between cycles. The
network returns to $V_{\mathrm{BUS}}$. Thus, \des{C}{HVCL} withstands only
$V_{\mathrm{cl}}$. The resistor also returns the discharged energy to the input
rail.

\subsubsection{Leakage energy and the sizing current}

Per cycle the leakage inductance delivers $E_{\mathrm{lk}} = \tfrac{1}{2}
L_{\mathrm{lk}} I_{\mathrm{pk}}^2$ into the clamp. Two currents matter:
\begin{itemize}[nosep]
  \item \textbf{Normal design-max operation}, $I_{\mathrm{pk}} =
    \SI{2.18}{\ampere}$ (\cref{sec:hv-op-converge}):
    $E_{\mathrm{lk}} = \tfrac{1}{2}\times\SI{0.15}{\micro\henry}\times 2.18^2
    \approx \SI{0.355}{\micro\joule}$, $E_{\mathrm{lk}} f_{\mathrm{sw}}
    \approx \SI{53.0}{\milli\watt}$. This sets the clamp's steady dissipation
    carried in the loss budget (\cref{sec:hv-op-losses}).
  \item \textbf{Worst case}: the current limit bounds the current at turn-off,
    so the sizing case is the guaranteed-maximum
    clamped peak including the CS filter's overshoot, \SI{2.91}{\ampere}
    (\cref{sec:hv-cs-sense}). Hiccup bounds how long the converter holds that
    peak (\cref{sec:hv-ss}) but not the peak itself.
    $E_{\mathrm{lk}} = \tfrac{1}{2}\times\SI{0.15}{\micro\henry}\times 2.91^2
    \approx \SI{0.635}{\micro\joule}$, $E_{\mathrm{lk}} f_{\mathrm{sw}}
    \approx \SI{94.9}{\milli\watt}$.
\end{itemize}
Clamp sizing controls the \des{Q}{HVSW} drain-voltage rating. Thus, it uses the
worst-case current. The loss budget uses the normal maximum-load current.

\subsubsection{Clamp energy multiplier}

The clamp absorbs more than the leakage energy. While \des{D}{HVCL} conducts,
the leakage current decreases linearly from $I_{\mathrm{pk}}$ to zero. The
driving voltage is $V_{\mathrm{cl}}-V_{\mathrm{refl}}$. Here,
$V_{\mathrm{refl}}=\SI{19.25}{\volt}$ (\cref{sec:hv-op-converge}):
\begin{align*}
  t_1 = \frac{L_{\mathrm{lk}} I_{\mathrm{pk}}}{V_{\mathrm{cl}} - V_{\mathrm{refl}}}.
\end{align*}
The delivered charge is $Q_{\mathrm{cl}}=\tfrac{1}{2}I_{\mathrm{pk}}t_1$.
It enters the clamp at $V_{\mathrm{cl}}$:
\begin{align*}
  E_{\mathrm{cl}} = V_{\mathrm{cl}} Q_{\mathrm{cl}}
    = \tfrac{1}{2} L_{\mathrm{lk}} I_{\mathrm{pk}}^2 \cdot
      \frac{V_{\mathrm{cl}}}{V_{\mathrm{cl}} - V_{\mathrm{refl}}}
    = E_{\mathrm{lk}} \cdot
      \frac{V_{\mathrm{cl}}}{V_{\mathrm{cl}} - V_{\mathrm{refl}}}.
\end{align*}
At steady state this equals \des{R}{HVCL}'s dissipation, giving the design
equation the rest of this section solves:
\begin{align*}
  \frac{V_{\mathrm{cl}}^2}{\des{R}{HVCL}}
    = E_{\mathrm{lk}} f_{\mathrm{sw}}
      \frac{V_{\mathrm{cl}}}{V_{\mathrm{cl}} - V_{\mathrm{refl}}}.
\end{align*}

\subsubsection{Target clamp voltage}

Clamp effectiveness sets the minimum $V_{\mathrm{cl}}$. Use at least twice the
reflected voltage (inferred). Switch stress sets the maximum. The total stress
includes $V_{\mathrm{BUS,max}}$, $V_{\mathrm{cl}}$, and the \des{D}{HVCL}
forward drop:
\begin{align*}
  V_{DS,\mathrm{pk}} = V_{\mathrm{BUS,max}} + V_{\mathrm{cl}} + V_{F(\des{D}{HVCL})}.
\end{align*}
This section limits semiconductor voltage stress to $80\,\%$ of the rating.
The analysis assumes a \SI{1.5}{\volt} \des{D}{HVCL} forward drop. This value
is near the high end for a small fast rectifier at ampere-level pulse current
(inferred). With $V_{\mathrm{BUS,max}} = \SI{12.6}{\volt}$:
\begin{align*}
  V_{\mathrm{cl}}^{\max} = (1-m)\times\SI{100}{\volt} - V_{\mathrm{BUS,max}}
    - V_{F(\des{D}{HVCL})} = 80 - 12.6 - 1.5 = \SI{65.9}{\volt}.
\end{align*}
This value is 3.4 times the reflected voltage. Therefore, switch stress sets the
limit. A higher clamp voltage decreases clamp loss and decreases switch margin.
The target is $V_{\mathrm{cl}} = \SI{65.9}{\volt}$.

\subsubsection{\des{R}{HVCL}}

Solving the energy-balance equation for \des{R}{HVCL} at the target, with the
worst-case-corner clamp power:
\begin{align*}
  \des{R}{HVCL} = \frac{V_{\mathrm{cl}}(V_{\mathrm{cl}} - V_{\mathrm{refl}})}
    {E_{\mathrm{lk}} f_{\mathrm{sw}}}
    = \frac{65.9 \times 46.65}{\SI{94.9}{\milli\watt}} \approx \SI{32.4}{\kilo\ohm}.
\end{align*}
Rounding down to the standard value $\des{R}{HVCL} = \SI{30}{\kilo\ohm}$
preserves margin; re-solving the same equation for the resulting clamp voltage
at each of the two operating points,
\begin{align*}
  V_{\mathrm{cl}}^2 - V_{\mathrm{refl}} V_{\mathrm{cl}}
    - E_{\mathrm{lk}} f_{\mathrm{sw}} \des{R}{HVCL} = 0
  \;\Longrightarrow\;
  V_{\mathrm{cl}}^{\mathrm{lim}} = \SI{63.8}{\volt}, \qquad
  V_{\mathrm{cl}}^{\mathrm{dm}} = \SI{50.7}{\volt}
\end{align*}
(sustained current-limit corner and normal design-max operation,
respectively). At the current-limit corner, switch stress is
$V_{DS,\mathrm{pk}} = 12.6 + 63.8 + 1.5 = \SI{77.9}{\volt}$ against the
\SI{100}{\volt} rating. The margin is \SI{22.1}{\percent}. During normal
operation, the clamp voltage is \SI{50.7}{\volt}. This is 2.6 times the
reflected voltage. The clamp dissipates \SI{85.5}{\milli\watt}
(\cref{sec:hv-op-losses}). At the current-limit corner, dissipation is
\SI{136}{\milli\watt}. Hiccup limits this condition to a
\SI{0.2}{\percent} duty cycle (\cref{sec:hv-ss}). Require a
\SI{0.18}{\watt} power rating and an \SI{80}{\volt} voltage rating.

\subsubsection{\des{C}{HVCL}}

The preceding calculations use a constant clamp voltage. Clamp ripple adds to
peak drain stress. Limit ripple to $2\,\%$ of $V_{\mathrm{cl}}$ at the
worst-case current. This limit keeps the stress error small.

The charge delivered during each cycle follows from
$E_{\mathrm{cl}}=V_{\mathrm{cl}}Q_{\mathrm{cl}}$:
\begin{align*}
  Q_{\mathrm{cl}} = \frac{E_{\mathrm{lk}}}{V_{\mathrm{cl}} - V_{\mathrm{refl}}}
    = \frac{\SI{0.635}{\micro\joule}}{44.55} \approx \SI{14.3}{\nano\coulomb}.
\end{align*}
At a $2\,\%$ ripple target ($\Delta V \le \SI{1.28}{\volt}$ at
$V_{\mathrm{cl}}^{\mathrm{lim}} = \SI{63.8}{\volt}$):
\begin{align*}
  \des{C}{HVCL}^{\min} = \frac{Q_{\mathrm{cl}}}{\Delta V}
    = \frac{\SI{14.3}{\nano\coulomb}}{\SI{1.28}{\volt}} \approx \SI{11.2}{\nano\farad}.
\end{align*}
Rounding up to $\des{C}{HVCL} = \SI{12}{\nano\farad}$ gives actual ripple
$Q_{\mathrm{cl}}/\des{C}{HVCL} \approx \SI{1.19}{\volt}$ ($1.9\,\%$) and a
discharge time constant $\tau = \des{R}{HVCL} \des{C}{HVCL} \approx
\SI{360}{\micro\second}$, about $54$ switching periods: the near-DC node the
treatment above assumes. Its voltage rating follows the corner clamp voltage,
$\ge\SI{80}{\volt}$.

\subsubsection{\des{D}{HVCL} requirements}

The design does not specify a part for \des{D}{HVCL}. It must meet these
requirements:
\begin{itemize}[nosep]
  \item \textbf{Reverse voltage.} It must block $V_{\mathrm{BUS,max}} +
    V_{\mathrm{cl}}^{\mathrm{lim}} = 12.6 + 63.8 = \SI{76.4}{\volt}$ during the
    on-time; at the same $20\,\%$ voltage-derating standard used throughout
    this section, that requires a rated reverse voltage of at least
    \SI{95.5}{\volt}, i.e.\ a $\ge\SI{100}{\volt}$-class part.
  \item \textbf{Forward current.} The leakage current decreases from
    \SI{2.91}{\ampere} to zero in approximately \SI{9.8}{\nano\second}. The
    charge is \SI{14.3}{\nano\coulomb}. Its average is approximately
    \SI{2.1}{\milli\ampere}. A rectifier that meets the voltage requirement will
    typically meet this short-pulse requirement (inferred).
  \item \textbf{Reverse recovery.} \des{D}{HVCL} must block again before the
    next turn-on, and the off-interval never falls below
    \SI{4.0}{\micro\second} (\cref{sec:hv-dcm}). A fast
    ($\le\SI{500}{\nano\second}$) or ultrafast ($\le\SI{75}{\nano\second}$) part
    clears that; a standard-recovery part is comparable to the window itself and
    is excluded.
\end{itemize}

\subsection{Switch (\des{Q}{HVSW}) stress and losses}\label{sec:hv-switch}

The clamp sets the maximum drain stress. The worst case is the sustained
current-limit condition (\cref{sec:hv-clamp}):
\begin{align*}
  V_{DS,\mathrm{pk}} &= V_{\mathrm{BUS,max}} + V_{\mathrm{cl}}^{\mathrm{lim}}
    + V_{F(\des{D}{HVCL})} = 12.6 + 63.8 + 1.5 = \SI{77.9}{\volt}
\end{align*}
This stress gives a $22.1\,\%$ margin below the IRL540N
\SI{100}{\volt} rating. It meets the $20\,\%$ derating target. In normal
maximum-load operation, the peak is approximately \SI{64.8}{\volt}. Peak current
spans \SIrange{2.2}{2.9}{\ampere}, below the \SI{36}{\ampere} rating. The
\SI{6.85}{\volt} gate supply fully enhances the logic-level MOSFET. It also
stays inside the $\pm\SI{16}{\volt}$ gate rating.

Total switch dissipation at design-max, at the typical turn-off overlap
($t_{\mathrm{ovl}}^{\mathrm{typ}}=\SI{36}{\nano\second}$, \cref{sec:hv-op-losses}):
conduction \SI{60}{\milli\watt} + turn-off \SI{140}{\milli\watt} + $C_{oss}$ few
\si{\milli\watt} $\approx \SI{0.20}{\watt}$, a rise of $\approx\SI{12}{\kelvin}$
on a free-standing vertical TO-220 at $\theta_{JA} \approx
\SI{62}{\kelvin\per\watt}$. At the maximum \SI{54}{\nano\second} turn-off
overlap, conduction loss is \SI{62}{\milli\watt}. Turn-off loss is
\SI{215}{\milli\watt}. Total dissipation is approximately
\SI{0.28}{\watt}, and the temperature rise is approximately
\SI{17}{\kelvin}.

\textbf{Turn-off overlap is the largest converter loss.}
\des{D}{HVGOFF} and \des{R}{HVGOFF} decrease gate-discharge resistance from
approximately \SI{11.5}{\ohm} to \SI{4.8}{\ohm}. They decrease overlap from
\SIrange{75}{114}{\nano\second} to
\SIrange{36}{54}{\nano\second}. \des{D}{HVGOFF} is reverse-biased during
turn-on. Thus, \des{R}{HVGON} controls turn-on.

\subsection{Output rectifier (\des{D}{HVRECT})}\label{sec:hv-rect}

Reverse stress occurs during the on-time:
\begin{align*}
  V_R &= V_{\mathrm{out}} + n V_{\mathrm{BUS,max}}
       = 191.5 + 10 \times 12.6 = \SI{317.5}{\volt}
\end{align*}
at the maximum output (\cref{sec:hv-achieved-range}). Secondary-side ringing
keeps the total below \SI{345}{\volt}. The UF4005 rating is
\SI{600}{\volt}, which gives at least \SI{42}{\percent} margin. Average forward
current is \SI{15}{\milli\ampere} (\cref{sec:hv-budget}). Peak forward current
is \SI{0.22}{\ampere}. Conduction loss is approximately
\SI{15}{\milli\watt}. In DCM, diode current reaches zero during the idle
interval. Thus, the switch does not hard-commutate the diode.

\subsection{Input and output capacitance}\label{sec:hv-caps}

\subsubsection{Input}

\des{C}{HVI1} is the \SI{100}{\micro\farad} polymer-aluminum bulk capacitor.
\des{C}{HVI2} is a \SI{4.7}{\micro\farad} ceramic capacitor. It closes the
high-frequency primary loop. \des{C}{HVBIAS} bypasses BIAS
(\cref{sec:hv-house-bias}). The primary RMS current is
\SI{1.04}{\ampere}. Its DC component is \SI{0.74}{\ampere}. Thus, the capacitor
bank carries \SI{0.73}{\ampere} RMS. \textbf{Use a polymer-aluminum bulk
capacitor.} A small general-purpose \SI{100}{\micro\farad} aluminum
electrolytic typically has a \SIrange{0.3}{0.5}{\ampere} ripple rating at
\SI{150}{\kilo\hertz} (inferred). This rating is insufficient. The calculated
bus ripple is
\begin{align*}
  \Delta V \approx
  \frac{(I_{\mathrm{pk}}/2 - I_{\mathrm{in}})\,t_{\mathrm{on}}}{C_{\mathrm{in}}}
  + I_{\mathrm{pk}}\,\mathrm{ESR}
  \approx \frac{\SI{1.57}{\micro\coulomb}}{\SI{105}{\micro\farad}} + 2.18\times0.04
  \approx \SI{102}{\milli\volt}\ \text{peak-to-peak},
\end{align*}
unobjectionable on a USB-PD input.

\subsubsection{Output}

$C_{\mathrm{out}} = \des{C}{HVO1}\,(\SI{4.7}{\micro\farad}$, radial polarized, polymer
aluminum$) + \des{C}{HVO2}\,(\SI{100}{\nano\farad}\ \text{film})$.

The tube ballasts make the load approximately constant current. The rail does
not need hold-up energy or a transient reservoir. It needs only
switching-ripple control. Capacitive ripple per cycle is
\begin{align*}
  \Delta V_{\mathrm{cap}} &=
  \frac{I_{\mathrm{out}}\,(T - t_{\mathrm{dem}})}{C_{\mathrm{out}}}
  = \frac{15.4\,\si{\milli\ampere} \times 5.6\,\si{\micro\second}}
         {\SI{4.8}{\micro\farad}}
  = \SI{18}{\milli\volt},
\end{align*}
exceeded by the ESR step $I_{s,\mathrm{pk}}\times\mathrm{ESR} \approx 0.22 \times
0.5 \approx \SI{108}{\milli\volt}$ peak-to-peak ($0.06\,\%$ of the rail) at
\des{C}{HVO1}'s $\approx\SI{0.5}{\ohm}$ ESR, with the film part shunting the
fast edges.
Through a \SI{10}{\kilo\ohm} ballast (\cref{sec:hv-load-digit}) this modulates
tube current by $\sim$\SI{11}{\micro\ampere} at \SI{150}{\kilo\hertz},
optically invisible. The small capacitance is sufficient, and beneficial: it
keeps the stored energy at \SI{69}{\milli\joule}, the start-up rise time at
\SI{17}{\milli\second} (\cref{sec:hv-ss}), and the bleed-down inside seconds
(\cref{sec:hv-bleeder}).

\textbf{\des{C}{HVO1} requires a \SI{250}{\volt} rating.} A
\SI{160}{\volt} part cannot withstand the \SI{191.5}{\volt} maximum setting.
The \SI{250}{\volt} rating gives at least $23\,\%$ margin. The selected part
must also have approximately \SI{0.5}{\ohm} ESR or less.

\subsection{Bleeder (\des{R}{HVBLEED}) and minimum load}\label{sec:hv-bleeder}

\des{R}{HVBLEED} discharges $C_{\mathrm{out}}$ after shutdown. It also supplies
the minimum load that prevents pulse skipping.

\subsubsection{Discharge}\label{sec:hv-bleeder-discharge}

Stored energy is modest, $\tfrac{1}{2} C_{\mathrm{out}} V^2 =
\tfrac{1}{2}\times\SI{4.8}{\micro\farad}\times 170^2 = \SI{69}{\milli\joule}$
(\cref{sec:hv-caps}), but the node must still bleed down after shutdown.

IEC 62368-1 classifies DC voltages at or below \SI{60}{\volt} as ES1. It
requires the stored voltage to reach this limit within \SI{2}{\second} after
disconnection. The tubes help discharge the rail. The digit branches pull the
rail to approximately \SI{148}{\volt} in \SI{28}{\milli\second}. They then stop
conducting. The colon branches continue to conduct toward their maximum
\SI{55}{\volt} maintaining voltage. Thus, the colon tubes can pull the rail
below the ES1 limit. \des{R}{HVBLEED} and the feedback divider set the discharge
time.

Two \SI{240}{\kilo\ohm} ballasts against $\des{R}{HVBLEED} = \SI{1.5}{\mega\ohm}$
in parallel with the feedback divider ($R_{\mathrm{div}} \approx
\SI{1.70}{\mega\ohm}$, \cref{sec:hv-divider}) give
\begin{align*}
  \des{R}{HVBLEED} \parallel R_{\mathrm{div}} &= 1.5\,\si{\mega\ohm} \parallel
    1.70\,\si{\mega\ohm} = \SI{797}{\kilo\ohm}, \\
  \tau &= \frac{C_{\mathrm{out}}}{2/\SI{240}{\kilo\ohm} + 1/\SI{797}{\kilo\ohm}}
       = \SI{0.50}{\second}, &
  V_\infty &= \frac{2\times\SI{55}{\volt}/\SI{240}{\kilo\ohm}}
                   {2/\SI{240}{\kilo\ohm} + 1/\SI{797}{\kilo\ohm}}
            = \SI{47.8}{\volt},\\
  t_{148\rightarrow 60\,\mathrm{V}} &= \tau
    \ln\frac{148 - 47.8}{60 - 47.8} = \SI{1.05}{\second}.
\end{align*}
The total time to \SI{60}{\volt} is approximately \SI{1.08}{\second}. This time
meets the \SI{2}{\second} limit with a \SI{46}{\percent} margin. It applies to a
board with all tubes installed. With the tubes removed, discharge takes
\SI{4.0}{\second}. This condition occurs during high-voltage adjustment.

The minimum-load requirement sets \des{R}{HVBLEED} to no more than
\SI{2.92}{\mega\ohm}. Use \SI{1.5}{\mega\ohm}. It dissipates
\SI{19}{\milli\watt} at the nominal output and \SI{24}{\milli\watt} at the
maximum output. The resistor sees \SI{191.5}{\volt}. With
\SI{20}{\percent} derating, require an element voltage rating of at least
\SI{239}{\volt}. Use a DIN0309 body.

\subsubsection{Minimum load}\label{sec:hv-bleeder-minload}

The controller cannot deliver less than one minimum-width pulse. At
\SI{12.6}{\volt} input and an approximately \SI{165}{\nano\second} minimum on
time:
\begin{align*}
  I_{\mathrm{pk}}^{\min} &= \frac{V_{\mathrm{BUS}} t_{\mathrm{on,min}}}{L_p}
    = \frac{12.6 \times 165\,\si{\nano\second}}{\SI{10}{\micro\henry}}
    = \SI{0.208}{\ampere},\\
  P_{\min} &= \tfrac{1}{2} L_p (I_{\mathrm{pk}}^{\min})^2 f_{\mathrm{sw}}\,\eta
    \approx \tfrac{1}{2}\times\SI{10}{\micro\henry}\times 0.0432 \times
      149.4\,\si{\kilo\hertz} \times 0.833
      \approx \SI{27}{\milli\watt}.
\end{align*}
Below this load, the output increases or the controller skips pulses. The
bleeder supplies \SI{19}{\milli\watt}. The divider supplies
\SI{17}{\milli\watt}. Their \SI{36}{\milli\watt} total is 1.34 times
$P_{\min}$. Thus, the rail regulates before the tubes ignite. The divider alone
is insufficient. The remaining load requires \des{R}{HVBLEED} to be no more
than \SI{2.92}{\mega\ohm}.

\subsection{Loop compensation}\label{sec:hv-comp}

\subsubsection{Plant}

A first-order model of a peak-current-mode DCM flyback is a controlled current
source into the output capacitance and the incremental load. At the nominal
point, the load conductance is
\begin{align*}
  G_L &= \frac{4}{\SI{7.5}{\kilo\ohm}}+
         \frac{2}{\SI{240}{\kilo\ohm}}+
         \frac{1}{\SI{1.5}{\mega\ohm}}+
         \frac{1}{\SI{1.70}{\mega\ohm}}
       = \SI{543}{\micro\siemens}.
\end{align*}
This conductance uses the slope of each load line. The tube maintaining voltages
and driver drop are constant offsets. They do not appear in the incremental
conductance. The output capacitance is $C_{\mathrm{out}} = \des{C}{HVO1} +
\des{C}{HVO2} = \SI{4.8}{\micro\farad}$ (\cref{sec:hv-caps}).

Modulator gain: the controller's PWM comparator sees COMP through its internal
attenuation $G_{\mathrm{COMP}} = 0.142$ and compares it against the sensed current
plus the fixed internal slope ramp of slope $S_e = \SI{40}{\milli\volt}\times
f_{\mathrm{sw}} = \SI{5975}{\volt\per\second}$ (\ds{lm5156h.pdf}{20}). Against
the sensed on-slope $S_n = R_S V_{\mathrm{BUS}}/L_p = 0.039\times9/\SI{10}{\micro\henry}
= \SI{35100}{\volt\per\second}$ (\cref{sec:hv-cs}):
\begin{align*}
  \frac{\partial I_{\mathrm{pk}}}{\partial v_c}
    &= \frac{G_{\mathrm{COMP}}}{R_S\,(1 + S_e/S_n)}
     = \frac{0.142}{0.039 \times 1.17}
     = \SI{3.11}{\ampere\per\volt}.
\end{align*}
Since $I_{\mathrm{out}} = \tfrac{1}{2} L_p I_{\mathrm{pk}}^2 f_{\mathrm{sw}}
\eta / V_{\mathrm{out}}$, the plant transconductance and its output resistance
follow:
\begin{align*}
  \frac{\partial I_{\mathrm{out}}}{\partial I_{\mathrm{pk}}}
    &= \frac{2 I_{\mathrm{out}}}{I_{\mathrm{pk}}}
     = \frac{2 \times 13.97\,\si{\milli\ampere}}{1.96}
     = 14.3\,\si{\milli\ampere\per\ampere},\\
     g_{ps} &= 14.3\,\si{\milli\ampere\per\ampere} \times \SI{3.11}{\ampere\per\volt}
     = \SI{44.3}{\milli\siemens} \text{ (per volt of COMP)},
\end{align*}
and because $I_{\mathrm{out}} \propto 1/V_{\mathrm{out}}$ at fixed peak current,
the converter adds conductance $I_{\mathrm{out}}/V_{\mathrm{out}} =
\SI{82.2}{\micro\siemens}$. The effective output resistance and dominant pole
are therefore
\begin{align*}
  R_{\mathrm{out}} &= \frac{1}{G_L+I_{\mathrm{out}}/V_{\mathrm{out}}}
       = \SI{1.60}{\kilo\ohm},\\
  G_0 &= g_{ps}R_{\mathrm{out}} = 70.9\ (\SI{37.0}{\decibel}), &
  f_p &= \frac{1}{2\pi R_{\mathrm{out}} C_{\mathrm{out}}}
       = \SI{20.7}{\hertz}.
\end{align*}
The DCM right-half-plane zero and the output capacitor's ESR zero
($\approx\SI{67}{\kilo\hertz}$ at \des{C}{HVO1}'s \SI{0.5}{\ohm}) both lie far
above the crossover derived below and are neglected.

\subsubsection{Compensator and loop}

The crossover $f_c$ must be above the \SI{20.7}{\hertz} plant pole. It must also
be below $f_{\mathrm{sw}}/10$, or approximately \SI{15}{\kilo\hertz}. The load
changes slowly because each tube has a dedicated driver. The USB-PD input is
also regulated. Therefore, the design uses a crossover of a few hundred hertz.
The Type-II zero is 0.7 to 0.8 times $f_c$. The high-frequency pole is more
than 10 times $f_c$.

The mid-band gain is $g_m\des{R}{HVCOMP}$. The error-amplifier transconductance
is \SI{2}{\milli\siemens} (\ds{lm5156h.pdf}{6}). The resistor and capacitor
ratios set the zero and high-frequency pole. Use
$\des{R}{HVCOMP}=\SI{10}{\kilo\ohm}$ in series with
$\des{C}{HVC1}=\SI{100}{\nano\farad}$. Shunt this network with
$\des{C}{HVC2}=\SI{220}{\pico\farad}$:
\begin{align*}
  f_z &= \frac{1}{2\pi \des{R}{HVCOMP} \des{C}{HVC1}} = \SI{159}{\hertz},
  &
  f_{hp} &= \frac{1}{2\pi \des{R}{HVCOMP}\,\frac{\des{C}{HVC1}\des{C}{HVC2}}{\des{C}{HVC1}+\des{C}{HVC2}}}
          = \SI{72.5}{\kilo\hertz},
\end{align*}
with integrator behavior below $f_z$ (the amplifier's
$g_m R_O = 2\,\si{\milli\siemens} \times 10\,\si{\mega\ohm} = \SI{86}{\decibel}$
DC gain, \ds{lm5156h.pdf}{25}, makes static regulation error negligible). The
divider contributes $H = V_{\mathrm{REF}}/V_{\mathrm{out}} = 1/170$.

Solving $|T(f)| = 1$ for the product of divider, compensator, and plant
(numerically, with all three break frequencies in place) confirms the target:
\begin{align*}
  f_c &\approx \SI{214}{\hertz} \approx \frac{f_{\mathrm{sw}}}{699},\\
  \varphi_m &= 180^\circ - \arctan\frac{f_c}{f_p} - 90^\circ
    + \arctan\frac{f_c}{f_z} - \arctan\frac{f_c}{f_{hp}}
    = 180 - 84.5 - 90 + 53.3 - 0.2 \approx 59^\circ.
\end{align*}
The zero is $0.74 f_c$. The high-frequency pole is 339 times the crossover
frequency. Both values meet their targets. The phase reaches
$-180^\circ$ only after the loop gain has decreased by tens of decibels. The
datasheet gives no tolerance for $G_{\mathrm{COMP}}$, $S_e$, or $g_m$. Verify
the loop with a transient test or Bode measurement.

\subsection{Soft start}\label{sec:hv-ss}

The datasheet soft-start equations apply only to boost and SEPIC converters.
The boost equation assumes an input-to-output conduction path. A flyback has no
such path. The \SI{10}{\micro\ampere} SS current charges \des{C}{HVSS}. The
error amplifier regulates FB to the lower of $V_{SS}$ and
$V_{\mathrm{REF}}$. Therefore,
$t_{SS} = \des{C}{HVSS} V_{\mathrm{REF}} / I_{SS}$.

The converter must supply enough current for the output to follow the ramp. At
the current limit, it charges $C_{\mathrm{out}}$ from zero to
$V_{\mathrm{out}}$. The tubes draw no current before ignition. The charge time
is
\begin{align*}
  t_{\mathrm{rise}} &\approx
  \frac{\tfrac{1}{2} C_{\mathrm{out}} V_{\mathrm{out}}^2}
       {\tfrac{1}{2} L_p I_{\mathrm{lim}}^2 f_{\mathrm{sw}}\,\eta}
  = \frac{\SI{69}{\milli\joule}}{\SI{4.1}{\watt}}
  \approx \SI{17.0}{\milli\second}.
\end{align*}
$\des{C}{HVSS}$ must be large enough that $t_{SS}$ exceeds $t_{\mathrm{rise}}$;
otherwise the output falls behind the commanded ramp and the error amplifier
saturates into the current limit. Setting $t_{SS} = t_{\mathrm{rise}}$ gives the
minimum:
\begin{align*}
  \des{C}{HVSS}^{\min} &= \frac{t_{\mathrm{rise}} I_{SS}}{V_{\mathrm{REF}}}
  = \frac{\SI{17.0}{\milli\second} \times \SI{10}{\micro\ampere}}{\SI{1.00}{\volt}}
  \approx \SI{170}{\nano\farad}.
\end{align*}
\paragraph{Hiccup constraint.}
The LM51561H specified here enables hiccup mode overload protection, which the
LM5156H option omits (\ds{lm5156h.pdf}{2}). A fault timer counts current-limited
cycles, but only once soft start has finished; 64 of them stop the converter and
pull SS down, it restarts after 32{,}768 cycles, and eight consecutive
unlimited cycles reset the count (\ds{lm5156h.pdf}{6,25}). At
\SI{149.4}{\kilo\hertz} that is a \SI{428}{\micro\second} fault window and a
\SI{219}{\milli\second} retry interval, a fault duty of $0.2\,\%$. The datasheet
states the requirement directly: soft start must be long enough not to trigger
hiccup once it finishes.

The converter charges $C_{\mathrm{out}}$ at its current limit. If the output
does not reach its target before soft start ends, current limiting continues.
Hiccup then operates after \SI{428}{\micro\second}. Therefore, $t_{SS}$ must
exceed the worst-case rise time. The calculation uses maximum output
capacitance, minimum current limit, and minimum efficiency. It gives
$t_{\mathrm{rise}}^{\max} \approx \SI{24.5}{\milli\second}$.

The nominal calculation would select \SI{220}{\nano\farad}. Its minimum
$t_{SS}$ is \SI{17.8}{\milli\second}, which is insufficient.
$\des{C}{HVSS} = \SI{470}{\nano\farad}$ gives a nominal
$t_{SS} = \SI{47.0}{\milli\second}$. Stacking
$\des{C}{HVSS}$'s own tolerance ($\pm10\,\%$, typical for this class of ceramic)
against $I_{SS}$ and $V_{\mathrm{REF}}$'s datasheet ranges
(\SIrange{9}{11}{\micro\ampere} and \SIrange{0.99}{1.01}{\volt},
\ds{lm5156h.pdf}{6}), worst-case-additive:
\begin{align*}
  t_{SS}^{\min} &= \frac{\SI{470}{\nano\farad}\times0.9 \times \SI{0.99}{\volt}}
  {\SI{11}{\micro\ampere}} \approx \SI{38.1}{\milli\second},
\end{align*}
\SI{13.5}{\milli\second} clear of $t_{\mathrm{rise}}^{\max}$. The full range is
\SIrange{38.1}{58.0}{\milli\second}.

Hiccup protection limits the duration of a short circuit. Without it, a shorted
tube socket can hold \des{Q}{HVSW} and \des{T}{HV} at the current limit.

\subsection{Summary}\label{sec:hv-summary}

The analysis uses the following inferred values: a
\SIrange{1.0}{2.5}{\volt} driver on-state drop, typical IRL540N
$R_{DS(\mathrm{on})}$ and $Q_G$, a \SI{0.5}{\volt} \des{D}{HVGOFF} forward
drop, a \SI{1.5}{\volt} \des{D}{HVCL} forward drop, a
\SI{50}{\milli\watt} transformer core-loss allowance, a
\SI{50}{\nano\second} comparator delay, and VCC regulator dropout below its
\SI{8}{\volt} specification point. The analysis does not quantify
\des{Q}{HVSW} output capacitance. The DCM idle margin must include its drain
ring-down interval.

\paragraph{Gate-charge effect on efficiency.}\label{sec:hv-summary-lever}
Bias current and turn-off overlap total
\SIrange{297}{400}{\milli\watt}. Both depend on \des{Q}{HVSW} gate charge. The
charge pump supplies this charge from a \SIrange{24}{32}{\volt} node. Thus,
gate-charge loss is approximately three times the loss from a direct input feed.
A \SI{20}{\nano\coulomb} part with similar on-resistance would increase maximum
load efficiency from \SI{83.3}{\percent} to approximately
\SI{87.1}{\percent}. No suitable through-hole part was identified (inferred).
The design therefore specifies the IRL540N.

{\footnotesize
\begin{longtable}{|p{0.23\textwidth}|p{0.335\textwidth}|p{0.335\textwidth}|}
\caption{170~V flyback converter: specification summary}\label{tab:hv-converter-summary}\\
\hline
\textbf{Parameter} & \textbf{Derived requirement} & \textbf{Specification / result} \\
\hline
\endfirsthead
\caption[]{170~V flyback converter: specification summary (continued)}\\
\hline
\textbf{Parameter} & \textbf{Derived requirement} & \textbf{Specification / result} \\
\hline
\endhead
\hline
\endlastfoot
$R_T$ & \SI{146.4}{\kilo\ohm} for \SI{150}{\kilo\hertz} & \SI{147}{\kilo\ohm}
  $\rightarrow$ \SI{149.4}{\kilo\hertz} \\
\hline
DA2032 frequency & characterized at \SI{100}{\kilo\hertz} only &
  \SI{149.4}{\kilo\hertz}; \SI{50}{\milli\watt} core-loss allowance \\
\hline
UVLO \des{R}{HVUV1}/\des{R}{HVUV2} & enable below the \SI{4.75}{\volt} minimum input & on
\SI{4.44}{\volt} / off \SI{3.97}{\volt}; \SI{63.4}{\kilo\ohm}/\SI{32.4}{\kilo\ohm} \\
\hline
Display inhibit \des{Q}{HVINH}/\des{Q}{HVREL} & default asserted with no
dependency on a rail outside the converter; released only on a contract whose
\emph{current} allowance carries the display (\cref{sec:pd-sequence}) & shunt across \des{R}{HVUV2}, gate
divided from $V_{\mathrm{BUS}}$ through \SI{100}{\kilo\ohm}/\SI{1}{\mega\ohm};
asserted by \SI{2.22}{\volt} in against a \SI{4.16}{\volt} earliest UVLO
release; \SI{199}{\micro\ampere}, $V_{GS(th)}\le\SI{2}{\volt}$ maximum \\
\hline
EN and UVLO & EN satisfied before UVLO releases & \SI{1.54}{\volt} $\ll$
\SI{4.44}{\volt} \\
\hline
PGOOD pull-up \des{R}{HVPG} & $\ge\SI{10}{\kilo\ohm}$, $\le\SI{18}{\volt}$,
  $\le\SI{1}{\milli\ampere}$ & \SI{100}{\kilo\ohm} to +3V3;
  \SI{33}{\micro\ampere} \\
\hline
VCC capacitor \des{C}{HVVCC} & \SIrange{1}{4.7}{\micro\farad} & \SI{2.2}{\micro\farad} \\
\hline
Bias charge pump \des{C}{HVPUMP}/\des{R}{HVPUMP} & hold VCC at \SI{6.85}{\volt} across the
whole input window; $V_{\mathrm{BIAS}}\ge\SI{8}{\volt}$ at the worst delivery corner &
\SI{4.7}{\nano\farad} C0G, \SI{82}{\ohm}; $V_{\mathrm{BIAS}}$ \SI{9.1}{\volt} worst
corner / \SI{20.4}{\volt} at \SI{12.6}{\volt} in; costs \SI{126}{\milli\watt} flat \\
\hline
\des{D}{HVSTART}, \des{D}{HVZ}, \des{C}{HVBIAS} & start-up path, fault clamp, reservoir &
\SI{40}{\volt} Schottky; \SI{30}{\volt} zener (non-conducting in normal operation);
\SI{1}{\micro\farad} \\
\hline
Tuning range & span \SIrange{170}{185}{\volt} (ignition margin) &
  \SIrange{164.8}{191.5}{\volt} nominal targets;
  \SIrange{146.5}{212.1}{\volt} informational stop-to-stop
  (\cref{tab:hv-ballast-divider}); end stops are impermissible setpoints \\
\hline
Digit ballast \des{R}{BALDIG} & rated digit current at every trim setpoint &
installed \SI{7.50}{\kilo\ohm}; permissible alternatives through
\SI{13.00}{\kilo\ohm}, with \des{R}{HVD3} matched per row
  (\cref{tab:hv-ballast-divider}) \\
\hline
Colon ballast \des{R}{BALCOL} & \SI{0.5}{\milli\ampere} brightness optimum, no
  rated band & \SI{240}{\kilo\ohm}; \SIrange{0.458}{0.569}{\milli\ampere}
  across the achieved range \\
\hline
\des{R}{HVD1}/\des{R}{HVD2} element stress & $<\SI{250}{\volt}$ per DIN0207 element &
  \SI{95.25}{\volt} max at design-max \\
\hline
$P_{\mathrm{out}}$ & Not applicable & \SI{2.38}{\watt} nominal / \SI{2.94}{\watt}
  design-max \\
\hline
Efficiency & Not applicable & design-max \SIrange{82.8}{83.6}{\percent} across the window,
peaking near \SI{7.4}{\volt}; nominal \SI{82.7}{\percent} / \SI{82.5}{\percent}
/ \SI{81.7}{\percent} at
\SI{4.75}{\volt} / \SI{9}{\volt} / \SI{12.6}{\volt} \\
\hline
Turn-off speedup \des{D}{HVGOFF}/\des{R}{HVGOFF} & reduce turn-off overlap; driver peak
  $\le\SI{1.5}{\ampere}$ & SB240S-E3/54 + \SI{3.3}{\ohm};
  $t_{\mathrm{ovl}}$ \SIrange{36}{54}{\nano\second}; \SI{1.32}{\ampere} peak \\
\hline
$I_{\mathrm{pk}}$ & Not applicable & \SI{1.96}{\ampere} nominal; \SI{2.18}{\ampere}
design-max; \SI{2.29}{\ampere} at $L_{\min}$ \\
\hline
DCM & $t_{\mathrm{on}}+t_{\mathrm{dem}} < T$ & $0.89\,T$ at the worst corner
(\SI{4.75}{\volt}, $L_p+10\,\%$, design-max): $10.5\,\%$ idle \\
\hline
Sense resistor \des{R}{HVCS} & \SIrange{37.8}{40.2}{\milli\ohm} worst-case-additive &
  \SI{39}{\milli\ohm} \\
\hline
Current-limit margin & Not applicable & $+1.7\,\%$ worst-case-additive; $+7.2\,\%$ RSS \\
\hline
CS filter \des{R}{HVCSF}/\des{C}{HVCSF} & $\tau$ \SIrange{30}{100}{\nano\second} & \SI{100}{\ohm},
  \SI{470}{\pico\farad}; $\tau=\SI{47}{\nano\second}$ \\
\hline
Clamped peak and saturation & Not applicable & \SI{2.91}{\ampere} and
  \SI{3.0}{\ampere} soft-saturation rating \\
\hline
Clamp \des{R}{HVCL}/\des{C}{HVCL} & $V_{\mathrm{cl}} \ge 2\times$ reflected; $20\,\%$ \des{Q}{HVSW}
  voltage derating at the current-limit corner & \SI{30}{\kilo\ohm},
  \SI{12}{\nano\farad}, both $\ge\SI{80}{\volt}$; $V_{\mathrm{cl}}$
  \SI{63.8}{\volt} corner / \SI{50.7}{\volt} design-max \\
\hline
\des{D}{HVCL} & $\ge\SI{100}{\volt}$; \SI{2.91}{\ampere}, \SI{14.3}{\nano\coulomb}
  pulse; fast or ultrafast recovery & Not applicable \\
\hline
\des{R}{HVCL} dissipation & rating $\ge 2\times$ sustained dissipation &
\SI{85.5}{\milli\watt} design-max; the \SI{136}{\milli\watt} corner is
hiccup-limited to a $0.2\,\%$ duty and does not set the rating. Rated
$\ge\SI{0.18}{\watt}$, $\ge\SI{80}{\volt}$ \\
\hline
\des{Q}{HVSW} $V_{DS}$ & $\le\SI{77.9}{\volt}$ at the current-limit corner &
  \SI{100}{\volt} rating; $22.1\,\%$ margin \\
\hline
\des{Q}{HVSW} dissipation & Not applicable & \SI{0.20}{\watt} typical / \SI{0.28}{\watt} maximum
$t_{\mathrm{ovl}}$; $\Delta T \approx \SIrange{12}{17}{\kelvin}$, TO-220, no
  heatsink \\
\hline
\des{D}{HVRECT} & $V_R < \SI{345}{\volt}$; \SI{15}{\milli\ampere} avg & UF4005:
  \SI{600}{\volt}, \SI{1}{\ampere} \\
\hline
Input bulk \des{C}{HVI1} & $\ge\SI{0.73}{\ampere}$ ripple capability &
\SI{100}{\micro\farad} polymer aluminum \\
\hline
Output caps \des{C}{HVO1}/\des{C}{HVO2} & rating $\ge\SI{250}{\volt}$ at the \SI{191.5}{\volt}
  envelope; switching-ripple control only & \SI{4.7}{\micro\farad} polymer +
  \SI{100}{\nano\farad} film, both \SI{250}{\volt} \\
\hline
Bleeder \des{R}{HVBLEED} & preload $\ge\SI{27}{\milli\watt}$ sets the value
  ($\le\SI{2.92}{\mega\ohm}$); $\le\SI{60}{\volt}$ within \SI{2}{\second}
  (IEC 62368-1 ES1) & \SI{1.5}{\mega\ohm}:
  \SI{36}{\milli\watt}, \SI{1.08}{\second} with tubes installed \\
\hline
Compensation \des{R}{HVCOMP}/\des{C}{HVC1}/\des{C}{HVC2} & $\varphi_m \ge 45^\circ$ & \SI{10}{\kilo\ohm},
  \SI{100}{\nano\farad}, \SI{220}{\pico\farad};
  $f_c\approx\SI{214}{\hertz}$, $\varphi_m\approx59^\circ$ \\
  \hline
  Soft start \des{C}{HVSS} & $t_{SS} > t_{\mathrm{rise}}^{\max} \approx
  \SI{24.5}{\milli\second}$, else hiccup trips at every power-up &
  \SI{470}{\nano\farad}; $t_{SS}$ \SI{47.0}{\milli\second} nominal,
  \SIrange{38.1}{58.0}{\milli\second} worst-case-additive \\
  \hline
  Hiccup (LM51561H) & bound the sustained current-limit corner & 64-cycle
  (\SI{428}{\micro\second}) fault window, \SI{219}{\milli\second} retry:
  $0.2\,\%$ fault duty \\
  \hline
  Spread spectrum & Not applicable & DITHOFF to VCC, dithering disabled: it would cost a
  third of the DCM idle margin (\cref{sec:hv-house-dithoff}) \\
  \hline
  Input current & carried by the negotiated contract (\cref{sec:hv-inhibit}) &
  \SI{744}{\milli\ampere} design-max at \SI{4.75}{\volt}; \SI{605}{\milli\ampere}
  nominal at \SI{4.75}{\volt} \\
  \hline
  \end{longtable}
  }

  Select the ballast and \des{R}{HVD3} as a pair from
  \cref{tab:hv-ballast-divider}. The gate speedup network decreases turn-off
  overlap from \SIrange{75}{114}{\nano\second} to
  \SIrange{36}{54}{\nano\second}. The bias pump is also necessary at minimum
  input. Without it, overlap reaches \SI{171}{\nano\second} and the converter
  cannot supply the maximum load. The minimum-input DCM idle margin is
  $10.5\,\%$. The additive current-limit margin is $1.7\,\%$. The bias pump
  dissipates \SI{126}{\milli\watt} because it supplies gate charge from the
  reflected-voltage node.

  Measure gate-discharge current, efficiency, current limit, and start-up
  voltage. These tests verify the main derived parameters.

  \subsection{Component selection}\label{sec:hv-parts}

  \Cref{tab:hv-parts} gives the selection criteria. Use through-hole parts when
  they are available. Use carbon-film bodies for the visible colon ballasts.
  The body column gives a requirement only when package size affects the
  analysis or an exact part is specified. Other parts can use any body that
  meets the electrical requirements.

  Resistor sizes use $50\,\%$ power derating and $20\,\%$ voltage derating. Use
  \SI{0.6}{\watt} and \SI{250}{\volt} for DIN0207 metal-film resistors
  (inferred). Use \SI{0.25}{\watt} and \SI{250}{\volt} for DIN0207 carbon-film
  resistors (inferred). Use a \SI{350}{\volt} DIN0309 voltage rating
  (inferred). Verify these values for the selected resistor series.

  {\footnotesize
  \begin{longtable}{|L{0.11\textwidth}|L{0.13\textwidth}|L{0.10\textwidth}|L{0.33\textwidth}|L{0.13\textwidth}|}
  \caption{170~V flyback converter: component selection}\label{tab:hv-parts}\\
  \hline
  \textbf{Designator} & \textbf{Class} & \textbf{Value} & \textbf{Required ratings and characteristics} & \textbf{Body} \\
  \hline
  \endfirsthead
  \caption[]{170~V flyback converter: component selection (continued)}\\
  \hline
  \textbf{Designator} & \textbf{Class} & \textbf{Value} & \textbf{Required ratings and characteristics} & \textbf{Body} \\
  \hline
  \endhead
  \hline
  \endlastfoot
  Not applicable & Flyback controller & LM51561H & hiccup-enabled variant (\cref{sec:hv-ss}) & HTSSOP-14 (PWP), thermal pad \\
  \hline
  \des{T}{HV} & Flyback transformer & DA2032-AL & 1:10; $L_p = \SI{10}{\micro\henry} \pm 10\,\%$; $I_{\mathrm{sat}} \geq \SI{3.0}{\ampere}$; $L_{\mathrm{lk}} \leq \SI{0.15}{\micro\henry}$; secondary DCR \SI{1.6}{\ohm} (\cref{sec:hv-op-losses}) & through-hole bobbin \\
  \hline
  \des{Q}{HVSW} & N-channel MOSFET & IRL540N & $V_{DS}\geq\SI{100}{\volt}$; $V_{GS}\geq\pm\SI{16}{\volt}$; logic-level, $R_{DS(\mathrm{on})}$ specified at $V_{GS}=\SI{5}{\volt}$ and $\leq\SI{55}{\milli\ohm}$ at \SI{6.85}{\volt}; $Q_{GD}\leq\SI{38}{\nano\coulomb}$; $I_D\geq\SI{3}{\ampere}$; $\theta_{JA}\leq\SI{62}{\kelvin\per\watt}$ (\cref{sec:hv-switch}) & TO-220, vertical \\
  \hline
  \des{Q}{HVINH} & N-channel MOSFET & TN2106N3-G & $V_{GS}\geq\pm\SI{12}{\volt}$; $V_{GS(th)}\leq\SI{2}{\volt}$ maximum (\cref{sec:hv-inhibit}) & TO-92 \\
  \hline
  \des{Q}{HVREL} & N-channel MOSFET & TN2106N3-G & $V_{DS}\geq\SI{12}{\volt}$; $V_{GS(th)}\leq\SI{2}{\volt}$ maximum (\cref{sec:hv-inhibit}) & TO-92 \\
  \hline
  \des{R}{BALDIG} & Fixed resistor & \SI{7.50}{\kilo\ohm} installed & 1\,\%; four; permissible alternatives require the matched \des{R}{HVD3} from \cref{tab:hv-ballast-divider}; $\geq\SI{0.32}{\watt}$ (\cref{sec:hv-load-digit}) & DIN0207 minimum \\
  \hline
  \des{R}{BALCOL} & Fixed resistor & \SI{240}{\kilo\ohm} & 5\,\%; two; $\geq\SI{0.16}{\watt}$; element voltage $\geq\SI{171}{\volt}$ (\cref{sec:hv-load-colon}) & DIN0207 carbon film \\
  \hline
  \des{R}{HVCS} & Current sense & \SI{39}{\milli\ohm} & 1\,\%; $\geq\SI{0.09}{\watt}$; low inductance (\cref{sec:hv-cs-sense}) & 2512 \\
  \hline
  \des{R}{HVCSF} & Fixed resistor & \SI{100}{\ohm} & 5\,\% (\cref{sec:hv-cs-filter}) & Not applicable \\
  \hline
  \des{C}{HVCSF} & Ceramic & \SI{470}{\pico\farad} & C0G; $\geq\SI{50}{\volt}$ (\cref{sec:hv-cs-filter}) & Not applicable \\
  \hline
  \des{R}{HVD1}, \des{R}{HVD2} & Fixed resistor & \SI{845}{\kilo\ohm} & 1\,\%; element voltage $\geq\SI{120}{\volt}$ (\cref{sec:hv-divider-stress}) & DIN0207 minimum \\
  \hline
  \des{R}{HVD3} & Fixed resistor & \SIrange{8.25}{9.09}{\kilo\ohm} & 1\,\%; matched to the \des{R}{BALDIG} row of \cref{tab:hv-ballast-divider} & Not applicable \\
  \hline
  \des{RV}{HVTRIM} & Cermet trimmer & \SI{2}{\kilo\ohm} & multiturn; $\pm10\,\%$ total resistance (\cref{sec:hv-divider-split}) & Not applicable \\
  \hline
  \des{R}{HVRT} & Fixed resistor & \SI{147}{\kilo\ohm} & 1\,\% (\cref{sec:hv-fsw}) & Not applicable \\
  \hline
  \des{R}{HVUV1} & Fixed resistor & \SI{63.4}{\kilo\ohm} & 1\,\% (\cref{sec:hv-uvlo}) & Not applicable \\
  \hline
  \des{R}{HVUV2} & Fixed resistor & \SI{32.4}{\kilo\ohm} & 1\,\% (\cref{sec:hv-uvlo}) & Not applicable \\
  \hline
  \des{R}{HVCOMP} & Fixed resistor & \SI{10}{\kilo\ohm} & 5\,\% (\cref{sec:hv-comp}) & Not applicable \\
  \hline
  \des{C}{HVC1} & Ceramic & \SI{100}{\nano\farad} & X7R; $\geq\SI{16}{\volt}$ (\cref{sec:hv-comp}) & Not applicable \\
  \hline
  \des{C}{HVC2} & Ceramic & \SI{220}{\pico\farad} & C0G; $\geq\SI{16}{\volt}$ (\cref{sec:hv-comp}) & Not applicable \\
  \hline
  \des{R}{HVGON} & Fixed resistor & \SI{10}{\ohm} & 5\,\% (\cref{sec:hv-switch}) & Not applicable \\
  \hline
  \des{R}{HVGOFF} & Fixed resistor & \SI{3.3}{\ohm} & 5\,\%; \SI{1.32}{\ampere} peak for \SI{36}{\nano\second} (\cref{sec:hv-switch}) & Not applicable \\
  \hline
  \des{D}{HVGOFF} & Schottky rectifier & SB240S-E3/54 & $V_R\geq\SI{20}{\volt}$; $V_F\leq\SI{0.5}{\volt}$ at \SI{1.3}{\ampere} (\cref{sec:hv-op-losses}) & DO-41 axial \\
  \hline
  \des{R}{HVCL} & Fixed resistor & \SI{30}{\kilo\ohm} & 5\,\%; $\geq\SI{0.18}{\watt}$; $\geq\SI{80}{\volt}$ (\cref{sec:hv-clamp}) & DIN0207 minimum \\
  \hline
  \des{C}{HVCL} & Metallized polypropylene film & \SI{12}{\nano\farad} & $\pm5\,\%$; $\geq\SI{80}{\volt}$; pulse-rated to \SI{300}{\ampere\per\micro\second}; minimum available lead spacing (\cref{sec:hv-clamp}) & Not applicable \\
  \hline
  \des{D}{HVCL} & Fast rectifier & To be selected & $V_R\geq\SI{100}{\volt}$; recovery $\leq\SI{500}{\nano\second}$; \SI{2.91}{\ampere} repetitive peak, \SI{14.3}{\nano\coulomb} per pulse (\cref{sec:hv-clamp}) & Not applicable \\
  \hline
  \des{D}{HVRECT} & Ultrafast rectifier & UF4005 & $V_R\geq\SI{600}{\volt}$; $I_{F(AV)}\geq\SI{1}{\ampere}$; $I_{F,\mathrm{pk}}\geq\SI{0.22}{\ampere}$ (\cref{sec:hv-rect}) & DO-41 axial \\
  \hline
  \des{C}{HVPUMP} & Ceramic & \SI{4.7}{\nano\farad} & C0G; $\pm5\,\%$; $\geq\SI{100}{\volt}$ (\cref{sec:hv-house-pump}) & Not applicable \\
  \hline
  \des{R}{HVPUMP} & Fixed resistor & \SI{82}{\ohm} & 5\,\%; $\geq\SI{0.21}{\watt}$; \SI{118}{\milli\ampere} peak (\cref{sec:hv-house-pump}) & DIN0207 minimum \\
  \hline
  \des{D}{HVPUMP}, \des{D}{HVRESET} & Fast rectifier & To be selected & $V_R\geq\SI{40}{\volt}$; recovery $\leq\SI{500}{\nano\second}$; \SI{118}{\milli\ampere} peak, \SI{6.6}{\milli\ampere} average (\cref{sec:hv-house-pump}) & Not applicable \\
  \hline
  \des{D}{HVSTART} & Schottky rectifier & To be selected & $V_R\geq\SI{40}{\volt}$; $V_F\leq\SI{0.4}{\volt}$ at \SI{6.6}{\milli\ampere} (\cref{sec:hv-house-pump}) & Not applicable \\
  \hline
  \des{D}{HVZ} & Zener & \SI{30}{\volt} & $\geq\SI{0.5}{\watt}$ (\cref{sec:hv-house-pump}) & Not applicable \\
  \hline
  \des{C}{HVBIAS} & Ceramic & \SI{1}{\micro\farad} & X7R; $\geq\SI{50}{\volt}$; $\geq\SI{0.5}{\micro\farad}$ effective at \SI{21}{\volt} bias (\cref{sec:hv-house-pump}) & Not applicable \\
  \hline
  \des{C}{HVVCC} & Ceramic & \SI{2.2}{\micro\farad} & X7R; \SIrange{1}{4.7}{\micro\farad}; $\geq\SI{16}{\volt}$ (\cref{sec:hv-house-vcc}) & Not applicable \\
  \hline
  \des{C}{HVSS} & Ceramic & \SI{470}{\nano\farad} & X7R; $\pm10\,\%$; $\geq\SI{16}{\volt}$ (\cref{sec:hv-ss}) & Not applicable \\
  \hline
  \des{C}{HVI1} & Polymer aluminum & \SI{100}{\micro\farad} & $\geq\SI{25}{\volt}$; ripple $\geq\SI{0.73}{\ampere}$ at \SI{150}{\kilo\hertz}; ESR $\leq\SI{50}{\milli\ohm}$ (\cref{sec:hv-caps}) & Not applicable \\
  \hline
  \des{C}{HVI2} & Ceramic & \SI{4.7}{\micro\farad} & X7R; $\geq\SI{25}{\volt}$ (\cref{sec:hv-caps}) & Not applicable \\
  \hline
  \des{C}{HVO1} & Polymer aluminum & \SI{4.7}{\micro\farad} & $\geq\SI{250}{\volt}$; ESR $\approx\SI{0.5}{\ohm}$ at \SI{100}{\kilo\hertz} (\cref{sec:hv-caps}) & Not applicable \\
  \hline
  \des{C}{HVO2} & Film & \SI{100}{\nano\farad} & $\geq\SI{250}{\volt}$ (\cref{sec:hv-caps}) & Not applicable \\
  \hline
  \des{R}{HVBLEED} & Fixed resistor & \SI{1.5}{\mega\ohm} & 5\,\%; element voltage $\geq\SI{239}{\volt}$ (\cref{sec:hv-bleeder-discharge}) & DIN0309 axial \\
  \hline
  \des{R}{HVPG} & Fixed resistor & \SI{100}{\kilo\ohm} & 5\,\% (\cref{sec:hv-house-pgood}) & Not applicable \\
\hline
\des{R}{HVINHPU} & Fixed resistor & \SI{100}{\kilo\ohm} & 5\,\% (\cref{sec:hv-inhibit}) & Not applicable \\
\hline
\des{R}{HVINHPD} & Fixed resistor & \SI{1}{\mega\ohm} & 5\,\% (\cref{sec:hv-inhibit}) & Not applicable \\
\hline
\des{R}{HVRELPD} & Fixed resistor & \SI{100}{\kilo\ohm} & 5\,\% (\cref{sec:hv-inhibit}) & Not applicable \\
\hline
\end{longtable}
}

Some parts require more than the nominal value. \des{D}{HVSTART} must be a
Schottky diode. A silicon diode can leave BIAS below its start-up minimum.
\des{C}{HVPUMP}, \des{C}{HVCSF}, and \des{C}{HVCL} require the specified
dielectric. Their capacitance sets charge or timing. \des{R}{HVBLEED} is a
safety-related part. A substitute must meet its element-voltage rating.

